\documentclass{article}

\usepackage[a-2u]{pdfx}
\usepackage[utf8]{inputenc}
\usepackage[T1]{fontenc}
\usepackage{lmodern}
\usepackage[czech]{babel}
\usepackage{array}
\usepackage{tabularx}
\usepackage{booktabs}
\usepackage{longtable}
\usepackage[ddmmyyyy]{datetime}
\usepackage{geometry}
\usepackage{enumitem}
\usepackage{url}
\usepackage[section]{placeins}
\emergencystretch=1.5em
\usepackage{xcolor}
\usepackage{amssymb}
\usepackage{listings}
\usepackage{tikz}
\usetikzlibrary{positioning,fit}

% --- Makety obrazovek: styly a pomocná makra převzatá ze specifikace ---
\tikzset{
  wf frame/.style={draw=black!70, rounded corners=3pt, fill=white},
  wf bar/.style={fill=gray!20},
  wf card/.style={draw=black!45, rounded corners=2pt, fill=gray!6},
  wf pre/.style={draw=black!25, fill=white},
  wf field/.style={draw=black!55, fill=white, rounded corners=1pt},
  wf popup/.style={draw=black!55, fill=white},
  wf h/.style={font=\linespread{1}\fontsize{8}{9}\selectfont\sffamily\bfseries, inner sep=1pt, anchor=west},
  wf t/.style={font=\linespread{1}\fontsize{7}{8}\selectfont\sffamily, inner sep=1pt, anchor=west},
  wf tt/.style={font=\linespread{1}\fontsize{6}{7}\selectfont\sffamily, inner sep=1pt, anchor=west},
  wf err/.style={wf tt, text=red!65!black},
  wf code/.style={font=\linespread{1}\fontsize{5.5}{6.6}\selectfont\ttfamily, inner sep=1pt, anchor=north west, align=left},
  wf tab/.style={font=\linespread{1}\fontsize{6}{7}\selectfont\sffamily, inner sep=1pt, anchor=north west},
  wf btn/.style={draw=black!55, rounded corners=2pt, fill=gray!22, font=\linespread{1}\fontsize{6}{7}\selectfont\sffamily\bfseries, inner sep=2.5pt, minimum height=4.4mm, anchor=west},
  wf primary/.style={wf btn, fill=blue!30},
  wf badge/.style={draw=black!35, rounded corners=4pt, fill=white, font=\linespread{1}\fontsize{6}{7}\selectfont\sffamily, inner sep=2pt, anchor=west},
}
\newcommand{\wfselect}[5]{%
  \draw[wf field] (#1,#2) rectangle (#3,#4);
  \node[wf t] at ({#1+0.12},{(#2+(#4))/2}) {#5};
  \fill[black!65] ({#3-0.34},{(#2+(#4))/2+0.06}) -- ({#3-0.12},{(#2+(#4))/2+0.06}) -- ({#3-0.23},{(#2+(#4))/2-0.08}) -- cycle;}
\newcommand{\wfinput}[5]{%
  \draw[wf field] (#1,#2) rectangle (#3,#4);
  \node[wf t] at ({#1+0.12},{(#2+(#4))/2}) {#5};}
\newcommand{\wfnav}[1]{%
  \fill[wf bar, rounded corners=3pt] (0,0) rectangle (16,-0.7);
  \fill[wf bar] (0,-0.35) rectangle (16,-0.7);
  \node[wf h] at (0.3,-0.35) {ORMConvertor};
  \node[wf t] at (9.4,-0.35) {\ifnum#1=1 \textbf{Translation}\else Translation\fi};
  \node[wf t] at (11.1,-0.35) {\ifnum#1=2 \textbf{Advisor}\else Advisor\fi};
  \node[wf t] at (12.4,-0.35) {Examples};
  \node[wf t] at (13.8,-0.35) {LDBC};
  \node[wf t] at (14.7,-0.35) {FIBEN};
  \node[wf t] at (15.5,-0.35) {SEDE};}

% --- Výpisy kódu ---
\definecolor{lstbg}{gray}{0.96}
\lstset{
  basicstyle=\linespread{1}\fontsize{7.5}{9}\selectfont\ttfamily,
  backgroundcolor=\color{lstbg},
  frame=single, framerule=0.2pt, rulecolor=\color{black!30},
  breaklines=true, breakatwhitespace=false,
  columns=fullflexible, keepspaces=true,
  showstringspaces=false, tabsize=4,
  xleftmargin=2pt, xrightmargin=2pt,
  aboveskip=4pt, belowskip=6pt,
  captionpos=b,
  literate={á}{{\'a}}1 {č}{{\v{c}}}1 {ď}{{\v{d}}}1 {é}{{\'e}}1 {ě}{{\v{e}}}1 {í}{{\'i}}1 {ň}{{\v{n}}}1 {ó}{{\'o}}1 {ř}{{\v{r}}}1 {š}{{\v{s}}}1 {ť}{{\v{t}}}1 {ú}{{\'u}}1 {ů}{{\r{u}}}1 {ý}{{\'y}}1 {ž}{{\v{z}}}1 {Á}{{\'A}}1 {Č}{{\v{C}}}1 {É}{{\'E}}1 {Ě}{{\v{E}}}1 {Í}{{\'I}}1 {Ř}{{\v{R}}}1 {Š}{{\v{S}}}1 {Ú}{{\'U}}1 {Ů}{{\r{U}}}1 {Ý}{{\'Y}}1 {Ž}{{\v{Z}}}1 {→}{{$\rightarrow$}}1 {×}{{$\times$}}1 {…}{{\dots}}1 {–}{{--}}1 {—}{{---}}1 {≥}{{$\geq$}}1 {✓}{{$\surd$}}1,
}
\renewcommand{\lstlistingname}{Výpis}

\renewcommand{\dateseparator}{.}
\renewcommand{\baselinestretch}{1.5}
\geometry{margin=2cm}
\setlength{\parskip}{3pt}

\newcolumntype{L}[1]{>{\raggedright\arraybackslash}p{#1}}
\newcolumntype{Y}{>{\raggedright\arraybackslash}X}

\def\ResearchProject{Research project}

% METADATA SECTION START
\def\StudentsFullName{Miroslav Štochel}
\def\StudyProgram{Informatika -- Softwarové a datové inženýrství}
\def\ProjectTitle{Rozšíření nástroje ORMorpher}
\def\ProjectType{\ResearchProject}
\def\SupervisorFullName{Ing. Pavel Koupil, Ph.D.}
\def\ConsultantsFullName{doc. RNDr. Irena Holubová, Ph.D., RNDr. Michal Kopecký, Ph.D.}
% METADATA SECTION END

\title{\ProjectTitle}
\author{\StudentsFullName}

\begin{document}
\shorthandoff{-}

\centerline{\Large \textbf{Report of research project}}
\centerline{ \textbf{Department of Software Engineering}}
\centerline{ \textbf{Faculty of Mathematics and Physics, Charles University}}
\bigskip
{\noindent\begin{tabular*}{\textwidth}{ >{\raggedleft}m{4cm} l}
 \textbf{Solver:} & \StudentsFullName \\
 \textbf{Study program:} & \StudyProgram \\[1em]
 \textbf{Project title:} & \ProjectTitle \\
 \textbf{Project type:} & \ProjectType \\[1em]
 \textbf{Supervisor:} & \SupervisorFullName \\
 \textbf{Consultants:} & \ConsultantsFullName \\[1em]
 \textbf{Repository:} & \url{https://github.com/MiraStochel/orm-convertor}, vydání \texttt{2.3.0} \\
 \textbf{Date:} & listopad 2026 \\
\end{tabular*}}

\bigskip
\noindent\textbf{Abstrakt.} Zpráva uzavírá výzkumný projekt navazující na schválenou specifikaci: nástroj ORMConvertor (v~publikacích ORMorpher) překládá entity, mapování a dotazy mezi šesti objektově-relačními frameworky dvou ekosystémů -- Dapper, NHibernate a EF~Core na .NET, Hibernate, EclipseLink a MyBatis na Javě -- přes společnou mezireprezentaci, doplňuje chybějící mapovací fakta z~katalogu databáze a ke každému převodu vrací strukturovanou diagnostiku. Dokument je zároveň provozní příručkou (co se staví a co běží, pro systém i pro testy a experimenty), uživatelským návodem (jak se nástroj používá, co a jak překládá, kde má meze a jak se projevují, se třemi propracovanými příklady), výběrem z~deníku vývoje (implementační zajímavosti a klíčová rozhodnutí) a zprávou o~experimentální části: matici překladů a metrikách korektnosti měřených experimentálním během, bězích nad katalogy LDBC, FIBEN a SEDE, hledání mezí nástroje a srovnání s~velkými jazykovými modely včetně lokálně nasazeného modelu Qwen3.5-9B. Reprodukce pro recenzenta je popsána v~samostatném souboru \texttt{docs/reproduction.md} v~repozitáři.

\tableofcontents
\clearpage

% ===== 01-uvod.local.tex =====
\section{Úvod}

Tato zpráva uzavírá výzkumný projekt \emph{Rozšíření nástroje ORMorpher} a navazuje na jeho schválenou specifikaci~\cite{spec}. Specifikace popsala výchozí stav prototypu, stav po dokončení .NET části (verze 1.2.0) a zbývající práci ve třech cílech: rozšíření .NET frameworků, javovské frameworky s~ověřením přes hranici ekosystémů a experimenty. Zpráva popisuje, co z~toho stojí ve vydání \texttt{2.3.0}, jak se nástroj nasazuje a používá, co se při vývoji ukázalo a co experimenty změřily. Analýzu a návrh, které specifikace obsahuje, neopakujeme; tam, kde se od ní výsledek liší, to říkáme výslovně (kap.~\ref{sec:odchylky}).

Repozitář je autoritativní zdroj: popis dnešního chování (\path{docs/architecture.md}), katalog toho, co se překládá a co ne (\texttt{docs/subset.md}), číslovaná rozhodnutí (\texttt{docs/decisions/}), záznamy běhů (\texttt{docs/runs/}) a hranice záruk (sekce \emph{Guarantees} kořenového \texttt{README.md}). Na rozhodnutí odkazujeme číslem (\emph{rozh.~113}), na záznam běhu datem; obojí je v~repozitáři dohledatelné jménem souboru.

\subsection{Motivace a smysl nástroje}

Databázové systémy mají nástroje pro migraci dat, evoluci schématu a propagaci změn do dotazů; aplikační kód nad databází takovou podporu nemá. Kdo mění objektově-relační framework -- protože se mění platforma, protože tým přechází z~.NET na Javu nebo naopak, protože se srovnává výkon kandidátů -- přepisuje entity, mapování a dotazy ručně a bez záruky, že nový dotaz vrací totéž co starý. ORMConvertor tuto práci dělá deterministicky: zdrojový parser přečte artefakty do frameworkově nezávislé mezireprezentace, nepovinná fáze doplní z~katalogu připojené databáze fakta, která zdroj neříká (typicky u~Dapperu a MyBatisu, kde entita nese jen jména a jazykové typy), a builder cílového frameworku z~ní napíše artefakty cíle. Co cíl nevyjádří, se nevynechá potichu: každý převod vrací záznamy o~ztrátách, konvencích, konfliktech, neúplnostech, doplněních z~katalogu, selháních a o~dotazech napsaných únikovou cestou nativním SQL, a knihu tvrzení, která u~každého mapovacího faktu říká, odkud přišel a zda do výstupu dorazil.

Smysl nástroje je dvojí. Pro migrujícího vývojáře je to pracovní pomůcka, která přeloží, co přeložit lze, a jmenovitě vypíše, co ne -- výstupem je výčet práce, která zbývá, ne iluze hotového projektu. Pro výzkum je to \emph{deterministická baseline}: přesně vymezená podmnožina konstrukcí, u~níž je každá hranice zapsaná rozhodnutím, a měřicí infrastruktura, která stejným metrem soudí artefakt nástroje i artefakt, který vyrobil člověk nebo jazykový model. Druhý smysl se v~průběhu projektu ukázal jako důležitější než první: srovnání s~jazykovými modely (kap.~\ref{sec:llm}) ukazuje, že silný model přeloží jednotlivé dotazy správně, ale nezná schéma databáze ani fakta připnuté verze frameworku, a tam, kde si je domyslí, soudce artefakt odmítne. Nástroj tato fakta bere z~katalogu a z~rozhodnutí podložených měřením.

\subsection{Vztah k~dalším projektům}

Projekt vznikl ve výzkumné skupině Adaptive Data Management (ADaM) na KSI MFF UK~\cite{adam-web}, jejíž práce se sbíhají ve frameworku \emph{MM-cat}~\cite{mmcat-intro}: platformně nezávislý konceptuální model nad logickými modely různých databázových systémů, založený na teorii kategorií~\cite{koupil-holubova-2022}, nad nímž se modelují multi-modelová data, dotazuje se nad nimi (MM-quecat, MM-evoque~\cite{mm-quecat,mm-evoque}), migrují se mezi modely a řídí se evoluce schématu (MM-evocat~\cite{mm-evocat}). Dokumentace frameworku popisuje tři směry, které se v~něm potkávají: jednotnou reprezentaci dat, dotazování nad ní a evoluci. Samostatný oddíl téže dokumentace, \emph{Database Mapping Optimization}~\cite{mmcat-aimm}, řeší volbu fyzické reprezentace dat jako optimalizační úlohu: nad pracovní zátěží měří latence v~PostgreSQL, MongoDB a Neo4j, učí z~nich modely latence a prohledává prostor mapování metodou Monte Carlo Tree Search; totéž téma nese nástroj mm-mapsearch~\cite{mm-mapsearch}.

ORMorpher stojí o~vrstvu výš. Zatímco MM-cat a Database Mapping Optimization volí reprezentaci \emph{dat} a databázový systém, ORMorpher překládá \emph{aplikační datovou vrstvu} -- entity, mapování a dotazy -- mezi frameworky, které nad zvolenou reprezentací stojí, a jeho Advisor volí framework (nebo kombinaci frameworků) pro zátěž dotazů z~naměřeného času a paměti. Oba směry se doplňují: jeden odpovídá na otázku, \emph{jak} mají být data uložena, druhý na otázku, \emph{čím} k~nim aplikace přistupuje, a oba ověřují volbu měřením, ne odhadem. Původní prototyp vytvořil Milan Abrahám v~diplomové práci~\cite{abraham-repo} a jeho návrh je popsán v~článku na ASE~2025~\cite{ormorpher-ase} a v~rozšířené časopisecké verzi~\cite{ormorpher-jss}; tento projekt pokračuje ve forku, bez sledování původního repozitáře.

Se souběžnými pracemi ve skupině se projekt dělí o~problém, ne o~kód: diplomová práce Martina Čorovčáka navrhuje univerzální objektový model a migraci aplikačního kódu mezi relačními, grafovými a dokumentovými systémy s~využitím velkých jazykových modelů; práce Ivony Oboňové migruje dotazy mezi relačními a grafovými systémy touž cestou. ORMorpher jde opačným směrem -- explicitními pravidly, heuristikami a mezireprezentací -- a právě proto dává pro LLM-based přístupy deterministickou baseline a společné měřidlo (kap.~\ref{sec:exp}).

\subsection{Návaznost na specifikaci}

Specifikace rozdělila práci do tří cílů a pěti milníků. Tabulka~\ref{tab:cile} říká, kde každý stojí; podrobnosti a přesné znění nároku nese kap.~\ref{sec:zaver}.

\begin{table}[ht]
\centering\small
\begin{tabularx}{\textwidth}{L{2.6cm} Y L{3.4cm}}
\toprule
\textbf{Cíl specifikace} & \textbf{Stav ve vydání 2.3.0} & \textbf{Požadavky} \\
\midrule
Cíl 1 -- .NET frameworky & hotový před specifikací (1.2.0); od té doby rozšířený o~výrazy, mezivýsledky, rekurzi, okenní funkce a nativní SQL jako únikovou cestu; dotazová matice ve všech 36 směrech šesti frameworků & F1--F6, F11, F14, S1--S3, S5--S7 \\
Cíl 2 -- javovské frameworky & hotový (2.0.0 a dál): Hibernate a EclipseLink jako dva profily nad sdílenou vrstvou Jakarta Persistence, MyBatis nad vrstvami jazyka Java a T-SQL, javovská testovací sada s~diferenčním ověřením, překlad přes hranici ekosystémů oběma směry & F7--F10, F12, F13 \\
Cíl 3 -- experimenty & T1--T3 nárokované v~užším rozsahu (vstup případové studie, matice kategorií, měřidlo experimentálního běhu); T4--T6 připravené protokolem a změřené na vzorcích, plné experimenty patří navazující diplomové práci; Advisor a T7 odložené za experimenty (rozh.~171) & T1--T7, F15 \\
\bottomrule
\end{tabularx}
\caption{Cíle specifikace a jejich stav.}
\label{tab:cile}
\end{table}

Harmonogram specifikace počítal s~odevzdáním v~listopadu; vydání \texttt{2.3.0} vzniklo 5.~listopadu~2026 nad záznamy soudce LDBC, běhu FIBEN a experimentálního běhu. Konec projektu podle schváleného záměru~\cite{zamer} je 1.~prosince~2026.

\subsection{Struktura dokumentu}

Kapitola~\ref{sec:nasazeni} je provozní dokumentace: co je v~repozitáři, co se staví do obrazů, co běží pro systém a co pro testy a experimenty, konfigurace a možné potíže při nasazení. Kapitola~\ref{sec:pouziti} je uživatelská dokumentace formou návodu: průchod překladovou obrazovkou, co nástroj jak překládá, kde má meze a jak se projevují, časté chyby ve vstupech a ve výpočtu, tři příklady od nejjednoduššího po složitý a návod k~experimentům. Kapitola~\ref{sec:denik} je výběr z~deníku vývoje: jak práce probíhala, implementační zajímavosti a rozhodnutí, co se nepovedlo napoprvé, kde se výsledek liší od specifikace a jak byl při práci využit asistent založený na AI. Kapitola~\ref{sec:exp} je experimentální část: metodika, výsledky nástroje, srovnání s~jazykovými modely a odkaz na reprodukční návod pro recenzenta. Kapitola~\ref{sec:zaver} shrnuje plnění požadavků a navazující práci.
% ===== 02-nasazeni.local.tex =====
\section{Nasazení a běh}\label{sec:nasazeni}

Tato kapitola říká, co repozitář obsahuje, co se z~něj staví, co pak běží při běžném provozu a co při testech a experimentech, a na které potíže lze při nasazení narazit. Kanonickou a průběžně udržovanou podobou téhož je provozní příručka \path{ORMConvertor/README.md} (anglicky); zde je stav k~vydání \texttt{2.3.0} a výklad, proč jsou věci uspořádané, jak jsou.

\subsection{Co je v~repozitáři a co se staví}

Repozitář má pět vrcholových složek. \texttt{ORMConvertor/} je nástroj sám -- řešení pro .NET~10 s~dvaadvaceti projekty, vedle něj javovská testovací sada \texttt{JavaTests/} (Maven, mimo \texttt{.sln}), importní konzolové projekty \texttt{FibenImport} a \texttt{SedeImport}, skripty databází v~\texttt{database/} a kontejnerová konfigurace. \texttt{docs/} je dokumentace (viz úvod), \texttt{benchmarks/} zděděné srovnání sedmi .NET ORM z~původního výzkumu, \texttt{diagrams/} diagramy prototypu a \texttt{notes/} zděděné výzkumné poznámky. Projekty řešení a jejich odpovědnost shrnuje tabulka~\ref{tab:projekty}; podrobný popis včetně závislostí nese \path{docs/architecture.md}, §2.

\begin{table}[ht]
\centering\small
\begin{tabularx}{\textwidth}{L{4.3cm} Y}
\toprule
\textbf{Projekt} & \textbf{Odpovědnost} \\
\midrule
\texttt{Model} & mezireprezentace: entity a mapování (\texttt{EntityMap}, seřazený složený klíč, vztahy s~páry sloupců), dotazové instrukce s~podmínkovým stromem \\
\texttt{AbstractWrappers} & rozhraní parserů a builderů, abstraktní buildery se šablonovou metodou, deskriptory cílových frameworků, typy diagnostických záznamů \\
\texttt{Common} & sdílené převodníky a jediný kompilační krok Roslynu \\
\texttt{DapperWrappers}, \texttt{EFCoreWrappers}, \texttt{NHibernateWrappers} & parser a builder každého .NET frameworku, každý ve vlastním projektu \\
\texttt{CSharpEntityParsing}, \texttt{LinqParsing}, \texttt{LinqBuilding} & sdílené čtení C\# entit, čtení a zápis LINQ \\
\texttt{JavaEntityParsing} & lexer a čtečka podmnožiny Javy, javovský typový slovník; žádný JVM \\
\texttt{JakartaPersistence} & sdílená vrstva JPA: anotace a \texttt{orm.xml}, zápis anotací, parser a zapisovač JPQL \\
\texttt{HibernateWrappers}, \texttt{EclipseLinkWrappers} & dva profily nad vrstvou JPA (odlišné výchozí hodnoty za týmiž anotacemi) \\
\texttt{MyBatisWrappers} & mapper nad vrstvami jazyka: Java pro doménovou třídu a rozhraní, T-SQL pro dokument mapperu \\
\texttt{TransactSql} & čtení a zápis T-SQL nad ScriptDom; používá ho Dapper, MyBatis, NHibernate i nativní SQL ostatních \\
\texttt{DatabaseCatalog} & jediné místo, které čte metadata databáze, a fáze doplnění \\
\texttt{OrmConvertor} & orchestrace: továrny, pipeline parser → doplnění → builder, \texttt{ConversionResult} \\
\texttt{SampleData} & vzorky, příklady, výňatky open-source projektů, katalogy LDBC, FIBEN a SEDE \\
\texttt{Advisor}, \texttt{AdvisorBenchmarking} & ILP model v~C nad GLPK přes P/Invoke; kompilace a měření generovaného kódu \\
\texttt{ORMConvertorAPI} & REST API v~ASP.NET Core a statický frontend ve \texttt{wwwroot} \\
\texttt{Tests} & sada xUnit~v3 po frameworcích a napříč; experimentální běhy, soudci katalogů, ovladače srovnání s~LLM \\
\bottomrule
\end{tabularx}
\caption{Projekty řešení a jejich odpovědnost.}
\label{tab:projekty}
\end{table}

\paragraph{Obrazy.} Vše, co běží v~kontejneru, vzniká ze dvou souborů \texttt{Dockerfile} a jednoho \texttt{docker-compose.yml} (rozh.~039: jeden soubor compose, dva profily). \path{ORMConvertorAPI/Dockerfile} má pět stupňů, tabulka~\ref{tab:stupne}; testovací stupně stojí před běhovým záměrně, aby prosté \texttt{docker build} dalo aplikaci. \texttt{database.Dockerfile} staví obraz SQL Serveru 2022, který při sestavení stáhne datové sady a při prvním startu svazku je načte: vzorovou databázi WideWorldImporters, \texttt{LdbcSnb} (LDBC Social Network Benchmark Interactive v1, scale factor 1, s~validační sadou, rozh.~110 a 117), \texttt{Fiben} (benchmark IBM Research, rozh.~149) a \texttt{Sede} (výpis dat dba.stackexchange.com pro dotazy Stack Exchange Data Exploreru, rozh.~208). Každý základní obraz je připnutý digestem vedle značky verze a balíčky se obnovují v~režimu \texttt{--locked-mode} z~commitnutých zámků (rozh.~186), takže dnešní sestavení dá tytéž bajty jako sestavení za rok.

\begin{table}[ht]
\centering\small
\begin{tabularx}{\textwidth}{L{2.6cm} L{4.6cm} Y}
\toprule
\textbf{Stupeň} & \textbf{Základ} & \textbf{Co dělá} \\
\midrule
\texttt{advisor-native} & \texttt{dotnet/sdk:10.0.402} & \texttt{gcc} a \texttt{libglpk-dev}; přeloží \texttt{Advisor/ilp.c} na \texttt{libadvisor.so}. Jediné místo, kde nativní knihovna Advisoru vzniká \\
\texttt{dotnet-build} & \texttt{dotnet/sdk:10.0.402} & zkopíruje každý \texttt{.csproj} s~jeho \texttt{packages.lock.json} jménem (kvůli cache restore), obnoví v~\texttt{--locked-mode}, publikuje API \\
\texttt{tests} & \texttt{dotnet-build} & sestavená sada xUnit; vstupní bod \path{dotnet test Tests/Tests.csproj --no-build} \\
\texttt{java-tests} & \path{maven:3.9.11-eclipse-temurin-25-noble} & javovská sada s~předem staženými závislostmi; vstupní bod \texttt{mvn -B -o test} (offline) \\
\texttt{runtime} & \path{dotnet/aspnet:10.0.13} + \texttt{libglpk40} & publikovaná aplikace a \texttt{libadvisor.so}; port 5072 \\
\bottomrule
\end{tabularx}
\caption{Stupně \texttt{ORMConvertorAPI/Dockerfile}.}
\label{tab:stupne}
\end{table}

\subsection{Spuštění systému}\label{sec:spusteni}

Pro běžný provoz stačí Docker; na hostiteli není třeba .NET SDK, JDK ani databáze. Z~adresáře \texttt{ORMConvertor/}:

\begin{lstlisting}[language=bash,caption={Spuštění systému.}]
git clone https://github.com/MiraStochel/orm-convertor.git
cd orm-convertor/ORMConvertor
docker compose up --build
\end{lstlisting}

Příkaz spustí dvě služby bez profilu, tabulka~\ref{tab:system}. První sestavení trvá podle připojení desítky minut: obraz databáze stahuje datové sady (LDBC SF~1 asi 222~MB, FIBEN 80~MB, výpis SEDE asi 200~MB) a obraz aplikace obnovuje balíčky. První start prázdného svazku načte WideWorldImporters ze zálohy, \texttt{LdbcSnb} asi za minutu, \texttt{Fiben} za dvě až tři minuty a \texttt{Sede} za několik minut, přičemž server už odpovídá; loader o~každé databázi zapíše rozšířenou vlastnost (\texttt{ldbc.dataset}, \texttt{fiben.dataset}, \ldots), podle níž při dalším startu pozná, že ji nemá načítat znovu. Aplikace čeká, až je databáze podle \emph{healthcheck} zdravá, a pak odpovídá na \url{http://localhost:5072/orm/}; Swagger UI je v~profilu \texttt{Development} na \texttt{/orm/swagger}, dokument OpenAPI vždy na \path{/orm/openapi/v1.json}.

\begin{table}[ht]
\centering\small
\begin{tabularx}{\textwidth}{L{2.6cm} Y L{3.6cm}}
\toprule
\textbf{Služba} & \textbf{Co běží} & \textbf{Port, procesory} \\
\midrule
\texttt{ormconvertor} & stupeň \texttt{runtime}: REST API a frontend pod \texttt{/orm}; proměnné \path{ConnectionStrings__CatalogDatabase} a \texttt{\_\_AdvisorDatabase} míří na \texttt{mssql\_db}; katalog je WideWorldImporters, nebo databáze z~\path{ORMCONVERTOR_CATALOG_DATABASE} (\texttt{LdbcSnb}, \texttt{Fiben}, \texttt{Sede}) & 5072; procesory 2--3 \\
\texttt{mssql\_db} & SQL Server 2022 s~WideWorldImporters, \texttt{LdbcSnb}, \texttt{Fiben} a \texttt{Sede} na trvalém svazku \texttt{mssql\_data}; \texttt{sa}/\texttt{Testingorms123} (jen pro vývoj) & 1444 na hostiteli; procesory 0--1 \\
\bottomrule
\end{tabularx}
\caption{Služby systému (\texttt{docker compose up --build}).}
\label{tab:system}
\end{table}

Ověření, že instance stojí: stránka \texttt{/orm/} se načte, \path{GET /orm/required-content} vrátí seznam jazyků, které každý zdroj čte, a překlad vzorku Dapperu do EF~Core na překladové obrazovce vrátí v~záznamu běhu stav katalogu \emph{reached} s~časem čtení. Zastavení je \texttt{docker compose stop} (data zůstávají), odstranění kontejnerů \texttt{docker compose down}, odstranění i svazku s~databázemi \texttt{docker compose down -v} -- další start pak načítá vše znovu.

\paragraph{Měřicí prostředí.} Compose připíná databáze na procesory 0--1 a každého klienta na 2--3, obraz databáze nastavuje \texttt{max server memory} 4096~MB a \texttt{MAXDOP}~1 a virtuální stroj Dockeru má mít čtyři procesory a 12~GB bez swapu (rozh.~145, 193). Tyto hodnoty jsou nastavené tak, aby měření Advisoru a běhů katalogů byla opakovatelná; pro pouhý provoz překladu jsou lhostejné, ale compose je nastavuje vždy -- s~méně než čtyřmi procesory Docker službu odmítne spustit, místo aby ji spustil na sdíleném jádru. Server má druhý režim, \emph{režim soudce} (rozh.~164, 166), do něhož ho přepne soubor \texttt{judge.env}: všechny procesory a jednovláknové plány pro souběžná čtení soudce LDBC. Dokud v~něm běží, musí každý příkaz compose nad ním nést \path{--env-file judge.env}, jinak compose databázi znovu vytvoří v~měřicím režimu a přeruší běžící přehrání (kap.~\ref{sec:problemy}).

\subsection{Běh bez Dockeru}

Překlad funguje kdekoli, kde je .NET~10 SDK:

\begin{lstlisting}[language=bash,caption={Spuštění z~.NET CLI.}]
cd ORMConvertor
dotnet run --configuration Release --launch-profile http --project ORMConvertorAPI/ORMConvertorAPI.csproj
\end{lstlisting}

Profil \texttt{http} nastaví prostředí \texttt{Development}, otevře \url{http://localhost:5072/orm/} a čte uživatelská tajemství (\emph{user secrets}), kde je u~nás \path{ConnectionStrings:CatalogDatabase}; bez připojení ke katalogu nástroj překládá dál, jen s~konvencemi místo doplnění a se stavem katalogu \emph{not configured} v~záznamu běhu. Ve Visual Studiu je spouštěcím projektem \texttt{ORMConvertorAPI}. Reálná instance se publikuje příkazem \texttt{dotnet publish} a spouští \path{dotnet ORMConvertorAPI.dll} \emph{z~adresáře publikace} -- pracovní adresář je kořen obsahu, kde aplikace hledá \texttt{wwwroot} a \texttt{appsettings.json}; spuštěná odjinud odpovídá na API, ale \texttt{/orm/} vrací prázdnou 404 (náprava: \texttt{--contentRoot}). Výchozí prostředí publikované instance je \texttt{Production}: bez Swaggeru a bez tajemství, konfigurace přichází proměnnými prostředí. Instance naslouchá bez TLS a bez přihlášení; veřejné nasazení patří za reverzní proxy, která ukončí HTTPS a drží prefix \texttt{/orm} (předpoklad důvěryhodné sítě, \path{docs/threat-model.md}).

Advisor mimo Docker nefunguje: potřebuje \texttt{libadvisor.so}, které vzniká jen ve stupni \texttt{advisor-native}, a SQL Server s~WideWorldImporters v~\path{ConnectionStrings:AdvisorDatabase}. Bez knihovny vrací jeho koncové body hlášku, ne pětistovku; překladu se to netýká.

\subsection{Konfigurace}

V~repozitáři nemá žádný připojovací řetězec hodnotu kromě jednorázových přihlašovacích údajů serverů, které si compose a CI samy vytvoří (S4). Klíče se píší s~\texttt{:} v~tajemstvích a \texttt{appsettings.json}, s~\texttt{\_\_} v~proměnných prostředí. Tabulka~\ref{tab:konfigurace} nese klíče provozu; proměnné jednotlivých běhů na vyžádání uvádí kap.~\ref{sec:behy} a úplný výčet \path{ORMConvertor/README.md}, \emph{Configuration}.

\begin{table}[ht]
\centering\small
\begin{tabularx}{\textwidth}{L{5.2cm} L{1.6cm} Y}
\toprule
\textbf{Klíč} & \textbf{Kdo} & \textbf{Význam; co se stane bez něj} \\
\midrule
\path{ConnectionStrings:CatalogDatabase} & API & doplnění z~katalogu. Bez něj jen konvence, hlášené stavem katalogu; F4 a F6 pak přes rozhraní neplatí. Nedostupný server se čeká 5~s bez opakování (rozh.~198), ale jen když zdroj dat jmenuje port nebo protokol \\
\path{ConnectionStrings:AdvisorDatabase} & API & měření Advisoru. Bez něj Advisor jmenuje klíč a skončí \\
\path{ConnectionStrings:TestDatabase} & .NET sada & databázové testy. Bez něj přeskočí s~důvodem; s~\path{ORMCONVERTOR_REQUIRE_TEST_DATABASE=1} selžou \\
\path{ORMCONVERTOR_TEST_JDBC_URL}, \path{ORMCONVERTOR_API_URL} & Java sada & JDBC URL s~uživatelem a heslem; běžící instance nástroje s~\texttt{/orm}, z~níž sada bere artefakty (rozh.~078) \\
\path{ASPNETCORE_ENVIRONMENT}, \texttt{ASPNETCORE\_URLS} & API & \texttt{Development} zapíná Swagger a tajemství; adresa naslouchání \\
\path{Parsing:MaxNestingDepth}, \path{Parsing:MaxQueryTextLength} & API & strop hloubky zanoření textu i modelu (výchozí 128) a délky textu jednoho dotazu (výchozí 1~MiB), rozh.~092, 191, 194; \texttt{0} vypíná. Hlubší nebo delší vstup je záznam \emph{Failure}; stropy jsou věcí provozovatele, ne požadavku \\
\path{ORMCONVERTOR_CATALOG_DATABASE}, \texttt{LDBC\_SCALE\_FACTOR}, \texttt{FIBEN\_DATASET}, \texttt{SEDE\_DATASET} & compose & katalog instance; scale factor LDBC nebo \texttt{none}; commit FIBEN nebo \texttt{none}; výpis SEDE nebo \texttt{none} \\
\bottomrule
\end{tabularx}
\caption{Klíče konfigurace provozu.}
\label{tab:konfigurace}
\end{table}

\subsection{Co běží pro testy}

Testovací sady jsou dvě a každou reprodukuje jeden příkaz na stroji, kde je jen Docker (S5); profil \texttt{test} compose souboru k~tomu staví vlastní prostředí, tabulka~\ref{tab:testprofil}.

\begin{lstlisting}[language=bash,caption={Obě sady v~kontejnerech.}]
docker compose --profile test build tests java_tests test_app test_app_ldbc test_app_cases test_app_dry test_db
docker compose --profile test run --rm tests         # .NET sada
docker compose --profile test run --rm java_tests    # javovská sada
\end{lstlisting}

\begin{table}[ht]
\centering\small
\begin{tabularx}{\textwidth}{L{3.1cm} Y}
\toprule
\textbf{Služba} & \textbf{Role} \\
\midrule
\texttt{test\_db} & obraz \texttt{mssql\_db} se scale factorem 0,1 a bez FIBEN a SEDE; první start načte WideWorldImporters, \texttt{LdbcSnb} a validační sadu do dvou minut \\
\texttt{test\_db\_init} & počká na zdravý server a na vlastnost \texttt{ldbc.validation}, vytvoří \texttt{ORMConvertorTests} a \texttt{ORMConvertorCaseStudies}; teprve po něm smějí sady startovat, aby nepřeskakovaly jen proto, že server ještě startoval \\
\texttt{tests} & stupeň \texttt{tests}; 6~GB, procesory 2--3; schémata si vytváří a maže sama (rozh.~016) \\
\texttt{test\_app}, \texttt{test\_app\_ldbc}, \texttt{test\_app\_cases}, \texttt{test\_app\_dry} & čtyři instance nástroje (stupeň \texttt{runtime}, 2~GB každá) s~katalogem nad testovací databází, nad \texttt{LdbcSnb}, nad databází případových studií a bez katalogu; javovská sada si z~nich bere artefakty přes HTTP \\
\texttt{java\_tests} & stupeň \texttt{java-tests}; 6~GB; běží po startu instancí \\
\bottomrule
\end{tabularx}
\caption{Služby profilu \texttt{test}.}
\label{tab:testprofil}
\end{table}

Každý \texttt{run} spustí SQL Server, počká na něj, vytvoří databáze a spustí celou sadu ve vlastním schématu; obě sady při tom přehrají i soudce LDBC nad prvními 1000 řádky validační sady SF~0,1 jako kontrolu. \textbf{Krok \texttt{build} není volitelný:} \texttt{run} sám znovu použije existující obraz a měří strom, z~něhož obraz vznikl -- zeleně a beze slova; jediným signálem je počet testů proti záznamu velikosti sady v~\texttt{docs/runs/README.md}. Na hostiteli s~.NET SDK a dosažitelným SQL Serverem běží .NET sada jako \path{dotnet test Tests/Tests.csproj --configuration Release}; testy, které databázi nemají, přeskočí s~důvodem. Javovská sada potřebuje JDK~25, Maven, databázi a běžící instanci, jejíž katalog míří na \emph{testovací} databázi (ne na vzorovou), jinak diferenční testy selžou hláškou o~nenalezené tabulce.

Průběžná integrace (\path{.github/workflows/ormconvertor-tests.yml}) běží při push a pull requestu do \texttt{main}, které mění \texttt{ORMConvertor/**} nebo workflow, ve třech úlohách: \texttt{test} (sestavení, \texttt{dotnet test} s~pokrytím proti SQL Serveru 2022 jako službě; databáze je vyžadovaná, takže přeskočený databázový test je selhání), \texttt{java-test} (Temurin 25; publikuje a spustí aplikaci, pak \texttt{mvn -B test}) a \texttt{dependencies} (známé zranitelnosti připnutých balíčků; jediná úloha spouštěná i týdně). CI nenačítá \texttt{LdbcSnb} ani \texttt{Fiben}, proto v~ní soudce LDBC a běhy katalogů přeskakují.

\subsection{Co běží pro experimenty}\label{sec:behy}

Experimenty nejsou součástí sestavení: běží \emph{na vyžádání}, za verzi nebo zásadní změnu, v~týchž kontejnerech profilu \texttt{test} nebo proti \texttt{mssql\_db} systému, a jejich čísla jsou záznam verze v~\texttt{docs/runs/}, ne nárok. Tabulka~\ref{tab:behy} vede druhy běhů; metodiku (co se před dlouhým během hlídá a v~jakém pořadí jde úplný běh) nese \texttt{docs/runs/README.md}, příkazy \path{ORMConvertor/README.md}, \emph{Tests}.

\begin{table}[ht]
\centering\small
\begin{tabularx}{\textwidth}{L{2.9cm} Y L{3.3cm}}
\toprule
\textbf{Běh} & \textbf{Co dělá; kde běží} & \textbf{Doba, záznam} \\
\midrule
soudce LDBC nad celou sadou & přehraje validační sadu LDBC Interactive v1 nad SF~1 (138~747 řádků): vkládání loaderem, každé čtení generovaným artefaktem každého frameworku, srovnání s~očekávanými řádky; \texttt{mssql\_db} v~režimu soudce, pro javovskou sadu instance nad \texttt{LdbcSnb} & asi 30~h na sadu; soubor \path{RRRR-MM-DD-ldbc-judge.md} \\
běh FIBEN, běh SEDE & každý dotaz katalogu do všech cílů s~katalogem, kompilace, srovnání řádků se zdrojem a se soudcem, medián času podle politiky rozh.~145; \texttt{mssql\_db} v~měřicím režimu, instance nad \texttt{Fiben} / \texttt{Sede} & FIBEN hodiny, SEDE přes 8~h na sadu; soubor \\
experimentální běh & tabulka T2 × T3: .NET půlka vypíše vstupy a soudí tři .NET cíle, javovská půlka tytéž vstupy do tří javovských cílů, souhrn počítá podíly; profil \texttt{test}, pro klíče \texttt{fiben:} navíc \texttt{mssql\_db} & asi hodina; soubor \\
běh mezí & syntetické řady podél jedné osy (joiny, podmínky, zanoření, poddotazy, mezivýsledky, vlastnosti, entity\ldots) od každého zdroje do každého cíle, proces na řadu, čas a paměť překladu; \texttt{test\_db} & .NET půlka 14~h, javovská 5~h; soubor \\
celý vstup & všechny vstupy výpisu sloučené do jednoho projektu na framework a šest projektů do šesti cílů, jen čas; \texttt{test\_db} a \texttt{mssql\_db} & minuty; soubor \\
cizí artefakty & adresář artefaktů s~manifestem (překlad od modelu, od člověka) projde stupni ověření v~procesu na artefakt s~mezí času (\texttt{foreign-run.sh}) & podle počtu; soubor \\
srovnání s~LLM & \texttt{cells.tsv} → prompty → pokusy (ručně ve webovém rozhraní, ovladačem nad API, nebo přes Claude Code) → import → experimentální běh nad vstupy vzorku → soud cizích artefaktů → souhrn & hodiny až den; soubor \\
\bottomrule
\end{tabularx}
\caption{Běhy na vyžádání.}
\label{tab:behy}
\end{table}

Každý dlouhý běh se spouští odpojeně, pojmenovaně a bez \texttt{--rm}, aby výsledek držel \texttt{docker logs} i poté, co shell, který ho spustil, skončil; soubory běhu (řádky po buňkách, souhrny, logy) zůstávají mimo git (rozh.~042) a do repozitáře jde jen záznam. Běh, který měří čas (FIBEN, SEDE, meze, celý vstup), chce na hostiteli vypnuté turbo a pozastavené aktualizace Windows (kap.~\ref{sec:problemy}); soudce LDBC a experimentální běh měří správnost a čas je jim lhostejný.

\subsection{Možné problémy při nasazení}\label{sec:problemy}

Následující potíže jsme potkali sami nebo je předvídáme; u~každé stojí projev a náprava.

\begin{enumerate}[leftmargin=*,itemsep=2pt]
\item \textbf{\texttt{run} bez \texttt{build}.} Sada běží nad stromem, z~něhož obraz vznikl, a je zelená. Projev: počet testů nesedí na záznam velikosti sady. Náprava: vždy dvojice \texttt{build} + \texttt{run}; u~dlouhého běhu je to nejdražší past, protože den měření stojí nad starým kódem.
\item \textbf{Paměť hostitele.} Systém (aplikace, SQL Server) a profil \texttt{test} (druhý SQL Server, čtyři instance po 2~GB, sada 6~GB) najednou přesáhnou 12~GB virtuálního stroje i 16~GB hostitele; proces hostitele je pak ukončen pro nedostatek paměti. Náprava: před během profilu \texttt{test} zastavit systém (\path{docker compose stop ormconvertor mssql_db}), odstranit zbylé kontejnery profilu, nespouštět hostitelský \texttt{dotnet test} vedle sestavení obrazu. Součet mezí všeho, co běží současně, se musí vejít do stroje (rozh.~193); experimentální běh proto pouští oba servery s~2048~MB.
\item \textbf{Virtuální stroj s~méně než čtyřmi procesory.} Compose připíná služby na procesory 0--1 a 2--3; s~menším strojem Docker službu odmítne. Náprava: \texttt{.wslconfig} (WSL2) se čtyřmi procesory a 12~GB, pak \texttt{wsl --shutdown} a start Docker Desktopu.
\item \textbf{Prázdná 404 na \texttt{/orm/} u~publikované instance.} Spuštěno z~jiného adresáře než z~kořene obsahu. Jediná stopa v~logu: \emph{Request reached the end of the middleware pipeline without being handled}. Náprava: \texttt{cd} do adresáře publikace, nebo \texttt{--contentRoot}.
\item \textbf{Katalog čeká desítky sekund.} Připojovací řetězec s~holým jménem hostitele nechá SqlClient hledat instanci přes SQL Browser; mez 5~s platí jen se jmenovaným portem nebo protokolem (\texttt{localhost,1444}, \texttt{tcp:host}). Náprava: port v~řetězci, případně vlastní \texttt{Connect Timeout}.
\item \textbf{Chybějící katalog u~zdroje Dapper nebo MyBatis.} Entita vyjde bez klíče (EF~Core jako \texttt{[Keyless]}), cíle, které klíč vyžadují (NHibernate, Hibernate, EclipseLink), ji odmítnou záznamem \emph{Failure}, a dotaz s~parametrem se odmítne záznamem, který jmenuje chybějící katalog. Není to vada: F4 a F6 bez databáze neplatí. Náprava: \path{ConnectionStrings:CatalogDatabase} na databázi se schématem vstupu.
\item \textbf{Katalog nad jinou databází, než jakou vstup popisuje.} Projev: záznam \emph{Incompleteness} \uv{No table matching the entity was found in the catalog}; u~javovské sady selhání diferenčních testů, protože instance překládala s~katalogem nad vzorovou databází místo testovací. Náprava: \path{ORMCONVERTOR_CATALOG_DATABASE} u~compose, \path{ConnectionStrings__CatalogDatabase} u~hostitelské instance.
\item \textbf{Port 1444 obsazený.} Zděděné benchmarky i \texttt{mssql\_db} publikují 1444; běží jen jedno z~nich.
\item \textbf{Režim soudce.} Dokud \texttt{mssql\_db} běží s~\texttt{judge.env}, musí ho nést každý příkaz compose, který se databáze dotkne; \texttt{up} závislé služby bez něj databázi znovu vytvoří v~měřicím režimu a přeruší přehrání. Skript, který dlouhý běh ukončuje, se smí databáze dotknout až poté, co sám přečetl kontejner běhu jako \texttt{exited} -- ukončení naplánované úlohy nechá \texttt{docker wait} naživu a navazující skript jinak databázi přetvoří pod běžícím soudcem (stalo se 24.~10.~2026 po 29~h přehrání, kap.~\ref{sec:nepovedlo}).
\item \textbf{Konce řádků.} \texttt{database/init-db.sh} a loadery běží v~linuxovém kontejneru; s~CRLF skončí \emph{bad interpreter}. \texttt{.gitattributes} drží \texttt{*.sh} na LF; po ručním zásahu ověřit \texttt{git ls-files --eol}.
\item \textbf{Nový projekt v~řešení.} Stupeň \texttt{dotnet-build} kopíruje každý \texttt{.csproj} a zámek jménem; projekt bez řádku v~tom seznamu rozbije sestavení obrazu. Hostitelský \texttt{dotnet restore} po posunu verze přepíše \texttt{packages.lock.json}, který je třeba commitnout, jinak CI a obraz v~\texttt{--locked-mode} selžou.
\item \textbf{Advisor mimo Linux.} Bez \texttt{libadvisor.so} vrací \texttt{/advisor-test} 400 s~chybou načtení knihovny; překlad běží. Windows sestavení knihovny neexistuje.
\item \textbf{Dlouhý běh přes noc.} Windows Update restartuje stroj kolem čtvrté ráno, po restartu nikdo není přihlášený a Docker Desktop neběží; uspání a hibernace běh přeruší také. Před během pozastavit aktualizace přes jeho očekávaný konec, vypnout uspání; po měřeném běhu vrátit stav procesoru (99~\% min.\ i max.\ pro měření) na výchozí.
\item \textbf{Databáze, která se nenačetla.} Loader skončí u~otázky, na kterou server neodpoví, a nechá databázi, jak byla; vlastnost \texttt{*.dataset} pak chybí nebo jmenuje jiný archiv a příští start načítá znovu. Projev v~logu \texttt{mssql\_db}; náprava podle hlášky loaderu (místo na disku, stažený archiv).
\item \textbf{Z~Git Bash.} \texttt{docker compose run} s~argumenty \texttt{-e} nesoucími \texttt{/app/\ldots} potřebuje \path{MSYS_NO_PATHCONV=1}, jinak shell přepíše cesty na windowsovské a soubory běhu přistanou v~adresáři toho jména.
\end{enumerate}
% ===== 03-pouziti.local.tex =====
\section{Použití}\label{sec:pouziti}

Kapitola je psaná jako návod: nejdřív co je vstupem a výstupem, pak průchod překladovou obrazovkou krok za krokem, co nástroj jak překládá, kde má meze a jak se meze projevují, časté chyby ve vstupech a ve výpočtu, tři propracované příklady a nakonec použití experimentů. Rozhraní je anglické, texty obrazovek proto citujeme v~originále.

\subsection{Vstup, výstup a co znamenají záznamy}

\textbf{Jednotka} je obsah jednoho zdrojového souboru v~jednom jazyce: třída C\# nebo Java (s~entitami, s~kódem, který předává dotazy, nebo s~obojím), mapovací dokument \texttt{hbm.xml} či \texttt{orm.xml}, mapper MyBatisu, skript SQL, holé HQL nebo JPQL. Jednotka deklaruje jen jazyk; co je v~ní entita a co předání dotazu, rozhodne zdrojový framework (rozh.~111). Převod je množina jednotek najednou, a proto jde celý projekt vložit po souborech nebo jako archiv ZIP, který rozbalí prohlížeč. Uživateli zůstává aplikační kód okolo -- volající metoda, transakce, DI, třída \texttt{DbContext} (její \texttt{DbSet} a \texttt{OnModelCreating} se čtou jako mapování, třída sama do výstupu nejde a převod to hlásí) a připojení do databáze, které do mezireprezentace nevstupuje (rozh.~029).

\textbf{Výstupem} jsou artefakty cíle po souborech: entita cíle (C\# nebo Java), mapovací dokument tam, kde ho cíl má (NHibernate \texttt{hbm.xml}, mapper MyBatisu), a za každý dotaz metoda v~kódu cíle, která dotaz předává jeho API, spolu s~holým textem dotazu (LINQ, HQL, JPQL, SQL). NHibernate dostává dotaz dvakrát: závazné HQL a vedle něj týž dotaz jako LINQ nad \texttt{session.Query<T>()} (rozh.~118). Nástroj nevydává spustitelný projekt (rozh.~040): soubor projektu, konfiguraci frameworku, jméno sestavení v~mapování NHibernatu a připojení doplní konzument. Ke každému převodu patří \textbf{záznam běhu}: identifikátor, verze nástroje a obou frameworků, stav katalogu a čas jeho čtení, deklarovaný dialekt (S6).

\textbf{Záznamy o~převodu} (F11) říkají, co se s~každým faktem stalo; tabulka~\ref{tab:zaznamy} nese jejich osm druhů v~pořadí, v~jakém je dobré je číst. Vedle záznamů stojí \textbf{kniha tvrzení} (\emph{Facts}): za každý mapovací fakt -- tabulka, schéma, sloupec, typ, délka, přesnost, nullabilita, klíč, strategie generování, sloupce cizího klíče, verze, unikátní omezení, nepersistovaná vlastnost -- řádek s~jeho hodnotou, původem (zdroj, katalog, odvozeno převodem), jednotkou, která ho vyslovila, a osudem (nesen, vypuštěn, negenerován), a kde to framework umí, s~pozicí ve vstupu i ve výstupu, kterou obrazovka podtrhne.

\begin{table}[ht]
\centering\small
\begin{tabularx}{\textwidth}{L{2.6cm} Y}
\toprule
\textbf{Druh} & \textbf{Význam a co s~ním dělat} \\
\midrule
\emph{Failure} & artefakt nevznikl: jednotka se nedá přečíst (s~řádkem a sloupcem), entita bez klíče míří do cíle, který klíč vyžaduje, dotaz by v~cíli vrátil jiné řádky (rozh.~053, 070), konstrukce mimo podmnožinu. Čte se první; zbytek převodu je hotový \\
\emph{Loss} & fakt zdroje, který cíl nevyjádří nebo mezireprezentace nenese (dědičnost, kaskády, nativní anotace, neunikátní index); výstup o~něm mlčí, záznam ho jmenuje \\
\emph{Convention} & výstup tvrdí něco, co zdroj neřekl -- výchozí generátor klíče, jméno klíčové třídy, jméno unikátního omezení; je to volba nástroje, ne fakt \\
\emph{Incompleteness} & chybí fakt, který cíl potřebuje: vztah bez sloupců, neznámý typ, tabulka nenalezená v~katalogu; cíl ho vynechá nebo nahradí konvencí \\
\emph{Supplied} & fakt dodal katalog databáze, ne zdroj (F4--F6); výstup o~něm říká víc než vstup \\
\emph{Conflict} & dva zdroje tvrdí totéž jinak (anotace vs.\ katalog, dva aliasy pro jednu vlastnost); platí priorita zdrojů (rozh.~017) a záznam říká, co prohrálo \\
\emph{Fallback} & jazyk cíle konstrukci nevysloví, dotaz vyšel celý nativním SQL deklarovaného dialektu přes API frameworku (rozh.~113); sémantika zachována, přenositelnost ne \\
\emph{Omitted} & NHibernate: tvar LINQ, který jeho provider nevysloví; závazné HQL stojí samo (rozh.~118) \\
\bottomrule
\end{tabularx}
\caption{Druhy diagnostických záznamů.}
\label{tab:zaznamy}
\end{table}

\subsection{Průchod překladovou obrazovkou}\label{sec:tutorial}

Rozcestník \texttt{/orm/} vede na překlad (\emph{Translation}), Advisor, výkladovou stránku s~živými příklady (\emph{Examples}) a na stránky katalogů LDBC, FIBEN a SEDE. Překladová obrazovka (obr.~\ref{fig:tutorial}) je jediná stránka s~oddíly nad sebou, takže se uživatel při opravě chyby vrací ke vstupu bez přepínání obrazovek (S7). Rozpracovaný vstup zůstává v~prohlížeči (\texttt{localStorage}), obnova stránky o~něj nepřijde (rozh.~056).

\begin{enumerate}[leftmargin=*,itemsep=2pt]
\item \textbf{Start from an example} (nepovinné). Výběr nabízí vzorek \emph{Customer} každého frameworku a popsané příklady z~výkladové stránky včetně výňatků skutečných projektů; \emph{Load example} naplní jednotky i směr. Pro první seznámení doporučujeme \emph{Dapper → EF Core: A source that states little}.
\item \textbf{Direction.} Řádek \emph{zdroj → cíl} se dvěma výběry a tlačítkem pro prohození; za ním volba dialektu doslovného SQL zdroje: nevysloveno (výchozí), SQL Server 2022, nebo jiný systém -- ten zastaví čtení doslovného SQL (rozh.~088), protože dotaz psaný pro jinou databázi by se nad SQL Serverem choval jinak.
\item \textbf{Input files.} \emph{Add files or a ZIP\ldots} přidá soubory nebo rozbalí archiv: jednotkou se stane každý soubor v~jazyce, který zdroj čte, jazyk se předvyplní podle přípony a ostatní soubory vyjmenuje sbalený seznam s~důvodem. \emph{Add empty unit} přidá prázdnou jednotku k~vepsání. Každá jednotka má jméno, jazyk a editor s~čísly řádků; \emph{Remove} ji odebere, \emph{Clear} vyprázdní vše. Nabídka jazyků je omezená na to, co zdroj čte (\path{GET /required-content}).
\item \textbf{Convert.} Validace na klientovi chytí prázdný seznam, prázdnou jednotku, jazyk mimo nabídku a nesprávně utvořené XML (s~řádkem) a ukáže chybu u~dotčené jednotky; syntaktické chyby C\#, Javy a SQL hlásí server záznamem \emph{Failure} s~řádkem a sloupcem, obrazovka řádek označí ve sloupci čísel a tlačítkem \emph{go to line} ho vybere. Po dobu volání je tlačítko zablokované a pod ním stojí stav; volání je jeden synchronní požadavek pod sekundu (S3).
\item \textbf{Result.} Hlavička nese záznam běhu. \emph{What you sent} a \emph{What came back} stojí vedle sebe a párují se podle entity, kterou server u~každého artefaktu vyslovil; u~panelu je \emph{Copy} a \emph{Download}, celek balí server do ZIP (\emph{Download all as ZIP}). Pás \emph{Facts} filtruje knihu tvrzení podle osudu a podtrhuje fakt ve vstupu i ve výstupu; pás \emph{Records} nese záznamy s~větou ke každému druhu a sloupci \emph{Unit}, \emph{Query} a \emph{Artifact}, záznamy \emph{Failure} první.
\end{enumerate}

\begin{figure}[htbp]
\centering
\begin{tikzpicture}[x=1cm,y=1cm]
\draw[wf frame] (0,0) rectangle (16,-11.2);
\wfnav{1}
\node[wf h, font=\linespread{1}\fontsize{11}{12}\selectfont\sffamily\bfseries] at (0.4,-1.1) {Translation};
\node[wf h] at (0.4,-1.65) {Start from an example};
\wfselect{0.4}{-1.9}{6.4}{-2.35}{Dapper → EF Core: A source that states little}
\node[wf btn] at (6.6,-2.125) {Load example};
\node[wf h] at (0.4,-2.8) {Direction};
\wfselect{0.4}{-3.05}{3.4}{-3.5}{Dapper 2.1.79}
\node[wf btn] at (3.55,-3.275) {$\rightleftarrows$};
\wfselect{4.1}{-3.05}{7.1}{-3.5}{Entity Framework Core 10.0.10}
\node[wf tt] at (7.4,-3.275) {Source SQL dialect:};
\wfselect{9.3}{-3.05}{11.9}{-3.5}{not stated}
\node[wf h] at (0.4,-3.95) {Input files};
\node[wf primary] at (0.4,-4.3) {Add files or a ZIP\dots};
\node[wf btn] at (2.65,-4.3) {Add empty unit};
\node[wf btn] at (4.55,-4.3) {Clear};
\draw[wf card] (0.4,-4.65) rectangle (7.9,-6.75);
\wfinput{0.55}{-4.78}{3.6}{-5.2}{Customer.cs}
\wfselect{3.75}{-4.78}{5.6}{-5.2}{C\#}
\node[wf btn] at (6.7,-4.99) {Remove};
\draw[wf pre] (0.55,-5.35) rectangle (7.75,-6.63);
\node[wf code] at (0.65,-5.42) {public class Customer\\\{\\~~~~public int CustomerID \{ get; set; \}\\~~~~public required string CustomerName \{ get; set; \}\\~~~~public DateTime AccountOpenedDate \{ get; set; \}\\~~~~public decimal? CreditLimit \{ get; set; \}\\\}};
\draw[wf card] (8.1,-4.65) rectangle (15.6,-6.75);
\wfinput{8.25}{-4.78}{11.3}{-5.2}{CustomerQuery.sql}
\wfselect{11.45}{-4.78}{13.3}{-5.2}{SQL}
\node[wf btn] at (14.4,-4.99) {Remove};
\draw[wf pre] (8.25,-5.35) rectangle (15.45,-6.63);
\node[wf code] at (8.35,-5.42) {SELECT c.CustomerName, c.CreditLimit\\FROM Sales.Customers AS c\\WHERE c.CreditLimit > 2000\\ORDER BY c.AccountOpenedDate DESC};
\node[wf h] at (0.4,-7.15) {Convert};
\node[wf primary] at (0.4,-7.5) {Convert};
\node[wf h] at (0.4,-8.0) {Result};
\draw[wf card] (0.4,-8.25) rectangle (15.6,-8.75);
\node[wf t] at (0.55,-8.5) {Dapper 2.1.79 $\to$ Entity Framework Core 10.0.10};
\node[wf badge, fill=green!12] at (6.3,-8.5) {catalog reached (572 ms)};
\node[wf t] at (9.6,-8.5) {ORMConvertor 2.3.0};
\node[wf t] at (12.3,-8.5) {run fa5731e7\dots};
\node[wf btn] at (0.4,-9.1) {Download all as ZIP};
\node[wf h] at (0.4,-9.55) {What you sent};
\node[wf h] at (8.1,-9.55) {What came back};
\draw[wf pre] (0.4,-9.75) rectangle (7.9,-10.4);
\node[wf code] at (0.5,-9.82) {Customer.cs ~~~~~~~~ CustomerQuery.sql};
\draw[wf pre] (8.1,-9.75) rectangle (15.6,-10.4);
\node[wf code] at (8.2,-9.82) {Customer.cs ~~ CustomerTransaction.cs ~~ CustomerQuery.cs};
\node[wf h] at (0.4,-10.6) {Facts};
\node[wf badge, fill=green!12] at (1.3,-10.6) {carried $\times$ 27};
\node[wf h] at (8.1,-10.6) {Records};
\node[wf badge, fill=green!12] at (9.4,-10.6) {Supplied $\times$ 27};
\node[wf badge, fill=yellow!20] at (11.3,-10.6) {Incompleteness $\times$ 4};
\node[wf badge, fill=yellow!20] at (13.7,-10.6) {Convention $\times$ 2};
\end{tikzpicture}
\caption{Překladová obrazovka ve vydání 2.3.0 (schematicky; hodnoty z~příkladu~1 v~kap.~\ref{sec:priklady}).}
\label{fig:tutorial}
\end{figure}

\paragraph{Ostatní stránky.} Výkladová stránka \emph{Examples} ukazuje potrubí převodu, výchozí vstup všech šesti frameworků a živé příklady po skupinách (uvnitř .NET, uvnitř Javy, přes hranici, celé domény, dotazy, aplikační kód, skutečné projekty); každý příklad se dá jedním tlačítkem načíst do překladu. Stránky \emph{LDBC}, \emph{FIBEN} a \emph{SEDE} ukazují katalogy referenčních dotazů se stavem každého dotazu, tabulku dotaz × cíl a převod vybraného dotazu nad běžící instancí; stránka FIBEN vysvětluje Kimovu typologii zanořených dotazů, stránka LDBC proč jsou dotazy na cesty rozepsané do řetězu úrovní. \emph{Advisor} přijímá entity, dotazy s~vahami a omezení (paměť, nejvyšší počet frameworků, kandidáti) a vrací srovnání všech změřených kandidátů s~vyznačeným přiřazením; měří jen Dapper a EF~Core a je vyňatý ze záruk, což stránka sama říká.

\paragraph{Přes API.} Totéž, co dělá obrazovka, dělá \texttt{POST /orm/convert} s~tělem \path{{sourceOrm, targetOrm, sources[{name, contentType, content}], declaredSourceDialect?}}; frameworky i jazyky jsou číselné výčty (\texttt{Dapper}~10, \texttt{NHibernate}~20, \texttt{EFCore}~30, \texttt{Hibernate}~40, \texttt{EclipseLink}~50, \texttt{MyBatis}~60; \texttt{CSharp}~90, \texttt{Java}~100, \texttt{XML}~30, \texttt{SqlQuery}~40, \texttt{HqlQuery}~50, \texttt{JpqlQuery}~80). Odpověď nese artefakty (\texttt{sources}, každý s~entitou a jednotkami), \texttt{records}, \texttt{facts} a záznam běhu. Chyba požadavku samého -- neznámý framework, chybějící pole -- je 400 s~\texttt{ProblemDetails}; vada nástroje při překladu není 400, ale záznam \emph{Failure} dotčené jednotky, entity nebo dotazu, a odpověď zůstává 200 s~tím, co se přeložilo (rozh.~197).

\subsection{Co jak překládá}\label{sec:co-jak}

Tabulka~\ref{tab:frameworky} říká, co každý framework čte jako zdroj a co píše jako cíl. Slovník dotazů je společný: projekce (i \texttt{DISTINCT} a konstanta pod aliasem), filtrace podmínkovým stromem s~AND/OR/NOT a závorkováním, joiny včetně levého, pravého, plného a křížového a joinu s~podmínkou nad rámec rovností, agregace, seskupení i podle výrazu, \texttt{HAVING}, řazení i podle agregátu, stránkování, poddotazy v~podmínce (\texttt{IN}, \texttt{EXISTS}, skalární porovnání, kvantifikované porovnání \texttt{ALL}/\texttt{ANY}), množinové operace, výraz všude, kde stojí operand (aritmetika, konkatenace, uzavřený slovník funkcí, \texttt{CASE}), parametry (pojmenované, poziční, zachycené hodnoty C\#, kolekce za \texttt{IN}), dotaz jako zdroj řádků (\texttt{WITH}, odvozená tabulka), rekurzivní \texttt{WITH} s~\texttt{OPTION (MAXRECURSION n)}, okenní funkce v~projekci a agregace do seznamu (\texttt{STRING\_AGG}). Každou konstrukci vymezuje rozhodnutí: nese se ve všech směrech, nebo je její odmítnutí, ztráta či zápis nativním SQL rozhodnuté a zapsané, ne zděděné z~toho, že na ni nedošlo.

\begin{table}[ht]
\centering\small
\begin{tabularx}{\textwidth}{L{1.9cm} Y Y Y}
\toprule
\textbf{Framework} & \textbf{Entity a mapování} & \textbf{Dotazy čte z} & \textbf{Dotazy píše jako} \\
\midrule
Dapper & třída C\# & SQL, SQL ve voláních Dapperu v~kódu & SQL s~metodou nad \texttt{IDbConnection} \\
EF Core & třída C\# s~anotacemi; fluent konfigurace v~\texttt{OnModelCreating} a \texttt{IEntityTypeConfiguration<T>} jako vymezená podmnožina & LINQ, nativní SQL v~kódu & LINQ \\
NHibernate & třída C\# + \texttt{hbm.xml} & LINQ, holé HQL, nativní SQL a HQL v~kódu & HQL a vedle něj LINQ nad \texttt{session.Query<T>()} \\
Hibernate & anotace \texttt{jakarta.persistence} (i~\texttt{orm.xml}), profil Hibernatu & JPQL a HQL Hibernatu, nativní SQL v~kódu & JPQL, HQL tam, kde JPQL tvar nemá (rekurze, limit rekurze hintem) \\
EclipseLink & totéž, profil EclipseLinku & JPQL, nativní SQL v~kódu & JPQL \\
MyBatis & doménová třída + \texttt{<resultMap>} & dokument mapperu (SQL), \texttt{@Select} na rozhraní & dokument mapperu a rozhraní \\
\bottomrule
\end{tabularx}
\caption{Co každý framework čte a píše.}
\label{tab:frameworky}
\end{table}

\textbf{Mapování} nese složené primární klíče jako seřazený seznam částí s~per-part strategií generování (ploché vykreslení, rozh.~006: \texttt{[PrimaryKey]} v~EF~Core, \texttt{<composite-id>} v~NHibernatu, \texttt{@IdClass} v~JPA), vícesloupcové cizí klíče s~uspořádanými páry sloupců, vztahy 1:1, 1:N a N:M -- ten jako explicitní spojovací entitu (rozh.~005) --, unikátní omezení, verzi a nepersistované vlastnosti. Hibernate a EclipseLink jsou dva profily nad jednou vrstvou: liší se výchozími hodnotami za týmiž anotacemi (\texttt{AUTO} je u~jednoho sekvence, u~druhého tabulka čítače; nationalizaci anotuje jen Hibernate), a převod proto čte \texttt{AUTO} profilem zdroje a píše do cíle mechanismus nativní dialektu, tedy \texttt{IDENTITY} (rozh.~187). MyBatis mapování nenese jako framework: \texttt{<resultMap>} je čten jako dvojice sloupec--vlastnost, \texttt{<id>} jako identita řádku výsledku, ne primární klíč.

\textbf{Úniková cesta.} Co jazyk cíle nevysloví, napíší EF~Core, NHibernate, Hibernate a EclipseLink celým dotazem v~nativním SQL deklarovaného dialektu přes API frameworku (\texttt{SqlQuery} nebo \texttt{FromSql}, \texttt{CreateSQLQuery}, \texttt{createNativeQuery}), vždy se záznamem \emph{Fallback} (rozh.~113): množinové operace a plný vnější join v~NHibernatu, mezivýsledek v~NHibernatu a EclipseLinku, rekurze v~EF~Core, NHibernatu a EclipseLinku, agregát nad celým výsledkem v~EF~Core, řetěz HQL nad 32 článků (rozh.~192). Dapper a MyBatis únikovou cestu nepotřebují, jejich jazykem SQL je.

\textbf{Doplnění z~katalogu} (F4--F6). Když zdroj neříká tabulku, schéma, typ, délku, nullabilitu, klíč, cizí klíč nebo unikátní omezení, převod si o~ně řekne katalogu připojené databáze dávkou za převod; co katalog dodá, je záznam \emph{Supplied}, co odporuje zdroji, \emph{Conflict} se zdrojovou hodnotou (rozh.~015, 017). Z~holé třídy Dapperu tak vznikne úplná entita EF~Core s~klíčem, typy, délkami a vztahy odvozenými z~cizích klíčů; dotaz, který porovnává parametr se sloupcem, si katalogem typuje parametr (rozh.~105).

\subsection{Limity a nemožnosti překladu}\label{sec:limity}

Hranice záruk je vyslovená na třech úrovních: \emph{vyňaté oblasti} (celé oblasti bez nároku), \emph{zúžení} (požadavek platí v~užším smyslu, než říkají jeho slova) a \emph{vyslovené meze} uvnitř nároku; všechny tři vede \texttt{docs/subset.md}, konstrukci po konstrukci jako \emph{konstrukce → proč → co nástroj udělá → kde rozhodnuto}. Čtyři vyňaté oblasti: (1)~Advisor a benchmarking -- bez testu funkce, měří dva frameworky ze šesti, kompiluje a spouští cizí kód bez izolace; (2)~dědičnost, komponenty a spojené tabulky -- NHibernate \texttt{<subclass>}, \texttt{<component>}, \texttt{<join>}, JPA \texttt{@Inheritance}, \texttt{@Embedded}, \texttt{@ElementCollection}, \texttt{@SecondaryTable}, EF~Core \texttt{[Owned]} a \texttt{[ComplexType]}: model nese jen plochou třídu, každý takový prvek dostane \emph{Loss} a vlastnost, kterou mapuje, z~modelu vypadne; (5)~jiný dialekt než SQL Server 2022 -- jediný, deklarovaný deskriptorem cíle; (6)~experimentální část za T1--T3, tedy LLM experimenty T4--T6 a T7 s~Advisorem. Čísla 3 a 4 patřila poddotazům a round-tripu NHibernatu, které z~výjimky vystoupily.

Tabulka~\ref{tab:limity} vede meze, na které uživatel narazí nejčastěji, a jak se projeví. Společné pravidlo nad všemi: \textbf{dotaz, který by v~cíli vrátil jinou množinu řádků, se nevydá} (rozh.~053, 065, 070) -- filtr, join, zdroj nebo seskupení, které parser nepřečte nebo které cíl věrně nevyjádří, odmítnou artefakt záznamem \emph{Failure}, zatímco projekce nebo řazení, které se liší jen v~tom, co se o~řádku říká, dostanou \emph{Loss}.

{\small
\begin{longtable}{L{3.6cm} L{9.6cm} L{3.2cm}}
\caption{Nejčastější meze a jejich projev.}\label{tab:limity}\\
\toprule
\textbf{Konstrukce} & \textbf{Proč a co nástroj udělá} & \textbf{Projev} \\
\midrule
\endfirsthead
\toprule
\textbf{Konstrukce} & \textbf{Proč a co nástroj udělá} & \textbf{Projev} \\
\midrule
\endhead
\bottomrule
\endlastfoot
dynamický příkaz MyBatisu (\texttt{<if>}, \texttt{<choose>}, \texttt{<bind>}) & rodina příkazů s~různými řádky, jeden člen ji nezastoupí; kanonický \texttt{<foreach>} je kolekční parametr a překládá se & \emph{Failure} se jménem značky \\
líné načtení v~EclipseLinku (\texttt{fetch = LAZY} na referenci) & bez weavingu tiše eager; model strategii načítání nenese & \emph{Loss}, artefakt \texttt{fetch} nevyslovuje \\
hodnota \texttt{CHAR(n)} v~EclipseLinku & EclipseLink 5.0.0 ve výchozí relaci odřízne koncové mezery, pět ostatních frameworků hodnotu předá doplněnou; měřítkem je hodnota SQL (rozh.~169) & do EclipseLinku konkatenace s~prázdným textem a konvertor doplnění; ze zdroje EclipseLink \texttt{TrimEnd}; vše \emph{Convention} \\
tvar LINQ, který provider NHibernatu nevysloví (pravý/plný join, agregát mimo \texttt{COUNT} nad \texttt{DISTINCT}, konkatenace nad nullable textem) & HQL je závazný tvar, LINQ doplněk & \emph{Omitted}, HQL stojí samo \\
zdroj Dapper nebo MyBatis bez katalogu & zdroj klíč nevysloví; NHibernate a JPA klíč vyžadují; parametr se bez katalogu netypuje & \emph{Failure} do NHibernatu, Hibernatu, EclipseLinku; EF~Core \texttt{[Keyless]}; dotaz s~parametrem odmítnut záznamem jmenujícím katalog \\
funkce mimo slovník (\texttt{CONVERT} s~délkou, \texttt{LEFT}, \texttt{SYSDATETIME()}, \texttt{ROUND} se třemi argumenty), agregát nad agregátem, výraz v~projekci bez aliasu & slovník funkcí je uzavřený a měřený v~obou ekosystémech (rozh.~071, 107) & \emph{Failure} v~podmínce, \emph{Loss} v~projekci \\
rámec okna, pořadová funkce bez řazení, \texttt{DISTINCT} v~okně, okno mimo projekci & cíle se liší; měřeno jen to, co je ve slovníku (rozh.~113, 161) & \emph{Loss} v~projekci, \emph{Failure} uvnitř mezivýsledku nebo mimo projekci \\
kolekční parametr v~nativním SQL EF~Core; LINQ skládané nad nativním SQL; dotaz skládaný za běhu (\texttt{QueryOver}, \texttt{CreateCriteria}) & API materializuje jinak nebo text neexistuje & \emph{Failure} \\
dotaz předaný jménem (\texttt{@NamedQuery}, \texttt{<query>}) & čte se tam, kde je definovaný, dotazovým průchodem jednotky, která ho nese (rozh.~143, 081) & záznam u~předání; dotaz vyjde z~definice \\
doplňkový balíček nad frameworkem (Dapper.Contrib, MyBatis-Plus) & tím frameworkem není; konvence balíčku platí jen v~projektu, který ho používá (rozh.~133) & \emph{Loss} se jménem balíčku a toho, co anotace tvrdila; tabulku a klíč doplní katalog \\
neznámý jazykový typ (\texttt{Instant}, mapa, výčet projektu, třída mimo převod) & slovník typů je uzavřený (rozh.~014, 075) & \emph{Incompleteness} s~doslovným typem; cíl v~druhém ekosystému ho nezkompiluje \\
zdroj s~jiným deklarovaným dialektem & doslovné SQL jiné databáze by nad SQL Serverem znamenalo něco jiného (rozh.~088) & \emph{Failure} dotazu, \emph{Loss} doslovného typu; katalog doplní s~jedním \emph{Conflict} \\
\end{longtable}}

Z~požadavků platí v~užším smyslu zejména: F7--F9 pro tři jmenované frameworky, ne pro javovský ekosystém; F11 syntaktickou správnost dokazují testy, ne \texttt{/convert}, protože překladová cesta cizí kód nekompiluje; F13 má měřítko porovnání pevné (desetinná místa u~dotazu, \texttt{NULL} vždy holé slovo), aby se shodly dvě sady ve dvou jazycích; F14 nenárokuje procházení repozitáře, jiný archiv než ZIP ani zobrazení mezireprezentace (rozh.~154 -- co z~ní uživatel potřebuje, říká kniha tvrzení); S7 má validaci na klientovi jako pomocníka, ne bránu. Úplné znění nese kořenové \texttt{README.md}, \emph{Guarantees}.

\subsection{Časté chyby ve vstupech}\label{sec:chyby-vstup}

\begin{itemize}[leftmargin=*,itemsep=2pt]
\item \textbf{Jednotka v~jiném jazyce, než deklaruje} (HQL v~jednotce označené SQL, \texttt{orm.xml} jako Java). Parser jazyka text nepřečte: \emph{Failure} jednotky s~řádkem a sloupcem, kde čtení skončilo. Náprava: jazyk z~nabídky podle přípony.
\item \textbf{Neúplná třída C\#} (chybějící středník, závorka). Roslyn by ji uměl \uv{dopsat}, převod ji ale odmítne na první chybě s~pozicí (rozh.~128), protože překlad z~dopsaného stromu by tvrdil něco, co ve zdroji není.
\item \textbf{Chybí entita, na kterou dotaz sahá.} Dotaz nad \texttt{Customer} bez jednotky s~třídou \texttt{Customer}: vztahy nemají cíl, cesty přes asociaci nejdou odvodit, dotaz je odmítnut nebo vyjde s~\emph{Incompleteness}. Vkládejte celý výňatek datové vrstvy, ne jeden soubor.
\item \textbf{Třída \texttt{DbContext}, služba, DTO.} Kontext se čte jako mapování a do výstupu nejde (záznam to říká); třída, která dotaz sama nepředává (služba volající repozitář, projekční DTO), se čte jako entita -- to je správně, pokud entitou je, a jinak ji z~převodu vynechte.
\item \textbf{Projekt bez katalogu u~Dapperu a MyBatisu.} Viz tab.~\ref{tab:limity}: bez databáze není odkud vzít klíč. Katalog musí být databáze se \emph{schématem vstupu}, ne ledajaká: záznam \emph{Incompleteness} \uv{No table matching the entity was found in the catalog} znamená, že instance míří na jinou databázi.
\item \textbf{Archiv s~cizími soubory.} ZIP projektu nese \texttt{.csproj}, \texttt{.json}, \texttt{.md}; obrazovka je vyjmenuje jako nepoužité, jednotkami jsou jen soubory v~jazycích zdroje. Položka archivu s~\texttt{..} v~cestě nebo s~nepovolenými znaky odmítne celý archiv s~400, které položku jmenuje (rozh.~199).
\item \textbf{Dvě jednotky deklarují touž entitu} (částečná třída, dvě verze souboru). Totožnost entity je dvojice jmenný prostor a název (rozh.~094): fakta se slučují podle priority zdrojů a rozpor je \emph{Conflict}; druhá definice je zpravidla omyl vstupu.
\item \textbf{Vstup z~jiného RDBMS.} MyBatis mapper psaný pro PostgreSQL nebo MySQL: gramatika T-SQL ho nepřečte (\emph{Failure} s~pozicí), a i když přečte, deklarujte dialekt \emph{another system}, aby se doslovné SQL nečetlo jako T-SQL.
\item \textbf{Příliš hluboký nebo dlouhý dotaz.} Strop zanoření (výchozí 128 úrovní závorek, \texttt{CASE}, prefixových operátorů nebo článků řetězu) a strop délky textu (1~MiB) odmítnou jednotku záznamem \emph{Failure}, který strop jmenuje; provozovatel je mění konfigurací, požadavek ne.
\end{itemize}

\subsection{Časté chyby ve výpočtu a jak se projevují}\label{sec:chyby-vypocet}

Nástroj nehádá. Tam, kde o~faktu nic neví, raději vydá méně a řekne to, než aby vydal víc a mlčel; projevy jsou vždy v~záznamech, nikdy v~tichém rozdílu výstupu.

\begin{itemize}[leftmargin=*,itemsep=2pt]
\item \textbf{Výstup říká víc než vstup.} Entita Dapperu bez atributů vyjde v~EF~Core s~\texttt{[Table]}, \texttt{[Column(TypeName)]}, \texttt{[MaxLength]}, \texttt{[Precision]} a \texttt{[Index(IsUnique)]}: každý z~těch faktů má záznam \emph{Supplied} a řádek knihy tvrzení s~původem \emph{catalog}. Bez katalogu by tytéž řádky nesly konvence cíle a záznam \emph{Convention}.
\item \textbf{Strategie generování klíče.} Katalog umí říct, že sloupec plní výchozí omezení, ne jaký mechanismus (rozh.~064): klíč zůstane \emph{Unspecified} s~\emph{Incompleteness} a cíl použije svou konvenci (\emph{Convention}): EF~Core generuje celočíselný klíč sám, NHibernate píše \texttt{assigned}, JPA \texttt{IDENTITY}.
\item \textbf{Vztah, který katalog zná a zdroj ne.} Cizí klíč v~katalogu bez navigační vlastnosti v~entitě vztah negeneruje (\emph{Incompleteness}); navigační vlastnost bez cizího klíče ve zdroji katalog doplní (\emph{Supplied}). Vlastnost typu entity bez jakéhokoli vztahu zůstane v~JPA \texttt{@Transient} s~\emph{Incompleteness} (rozh.~178).
\item \textbf{Dotaz vyšel v~nativním SQL.} Záznam \emph{Fallback} jmenuje konstrukci a dialekt; artefakt běží a vrací tytéž řádky (měřeno na čtvrtém stupni), ale není přenositelný mimo SQL Server. Kdo chce JPQL nebo LINQ, musí dotaz přepsat sám.
\item \textbf{Ztracené inicializátory a modifikátory přes hranici ekosystémů.} \texttt{= new ArrayList<>()} a \texttt{= []} nejsou literály, které oba jazyky píší stejně, proto se vypustí s~\emph{Loss}; \texttt{virtual}, \texttt{required}, \texttt{final} nemají protějšek a odpadnou (rozh.~076, 077).
\item \textbf{Verze.} \texttt{@Version int} Hibernatu do EF~Core vyjde jako \texttt{[ConcurrencyCheck]} s~\emph{Loss}: anotace EF~Core inkrementaci nevyjádří, zvyšovat ji musí aplikace. \texttt{[Timestamp] byte[]} do EclipseLinku vyjde jako sloupec jen pro čtení bez \texttt{@Version}, protože EclipseLink 5.0.0 verzi nad binárním sloupcem odmítá (rozh.~179; změřeno).
\item \textbf{Ekvivalence, která není totožnost.} \texttt{SUM} jazyka SQL nad prázdnou skupinou je \texttt{NULL}, \texttt{Sum} EF~Core je \texttt{COALESCE(SUM(x), 0)}; builder píše tvar se stejným výsledkem (rozh.~155), a proto vypadá LINQ pro \texttt{sum} složitěji, než by člověk napsal (viz příklad~3). Podobně \texttt{BETWEEN} vychází jako dvojice porovnání (pravidlo Q14 článku) s~\emph{Convention}.
\item \textbf{Vada nástroje.} Výjimka při čtení jednotky, generování entity nebo dotazu je \emph{Failure} té jednotky, entity či dotazu a řádek \texttt{Error} v~logu instance s~identifikátorem běhu (rozh.~197); zbytek převodu doběhne. Chybu hlaste s~identifikátorem běhu.
\end{itemize}
% ===== 03b-priklady.local.tex =====
\subsection{Tři příklady}\label{sec:priklady}

Všechny tři příklady jsou v~nástroji k~načtení volbou \emph{Start from an example}; výstupy níže jsou skutečné artefakty vydání 2.3.0 nad instancí s~katalogem (WideWorldImporters), zkrácené o~gettery a settery a o~opakující se záznamy.

\subsubsection*{Příklad 1 (základní): Dapper → EF Core, zdroj, který říká málo}

Vstupem jsou dvě holé třídy Dapperu bez jediného atributu a jeden dotaz v~SQL (výpis~\ref{lst:ex1-in}). Zdroj neříká tabulku, schéma, typy sloupců, klíč ani vztah; všechno dodá katalog.

\begin{lstlisting}[language={[Sharp]C},caption={Příklad 1, vstup: \texttt{Customer.cs}, \texttt{CustomerTransaction.cs} (zkráceno) a \texttt{CustomerQuery.sql}.},label={lst:ex1-in}]
namespace DapperEntities;
public class Customer
{
    public int CustomerID { get; set; }
    public required string CustomerName { get; set; }
    public DateTime AccountOpenedDate { get; set; }
    public decimal? CreditLimit { get; set; }
    public List<CustomerTransaction> Transactions { get; set; } = [];
}
public class CustomerTransaction
{
    public int CustomerTransactionID { get; set; }
    public int CustomerID { get; set; }
    public DateTime TransactionDate { get; set; }
    public decimal TransactionAmount { get; set; }
}
-- CustomerQuery.sql
SELECT c.CustomerName, c.CreditLimit
FROM Sales.Customers AS c
WHERE c.CreditLimit > 2000
ORDER BY c.AccountOpenedDate DESC
\end{lstlisting}

Výstup (výpis~\ref{lst:ex1-out}) je entita EF~Core s~tabulkou, schématem, databázovými typy, délkou, přesností, klíčem a unikátním omezením a dotaz v~LINQ se stejnou projekcí, filtrem a řazením. Záznam běhu říká \emph{catalog reached}, 572~ms čtení katalogu, verze Dapper 2.1.79 → EF~Core 10.0.10, nástroj 2.3.0.

\begin{lstlisting}[language={[Sharp]C},caption={Příklad 1, výstup: \texttt{Customer.cs} a dotaz.},label={lst:ex1-out}]
[Table("Customers", Schema = "Sales")]
[Index(nameof(CustomerName), IsUnique = true, Name = "UQ_Sales_Customers_CustomerName")]
public class Customer
{
    [Key]
    [Column(TypeName="int")]
    public required int CustomerID { get; set; }

    [Column(TypeName="nvarchar")]
    [MaxLength(100)]
    public required string CustomerName { get; set; }

    [Column(TypeName="date")]
    public required DateTime AccountOpenedDate { get; set; }

    [Column(TypeName="decimal")]
    [Precision(18, 2)]
    public decimal? CreditLimit { get; set; }

    public List<CustomerTransaction> Transactions { get; set; } = [];
}

public static IQueryable Query(DbContext ctx)
{
    return ctx.Set<Customer>()
        .Where(c => c.CreditLimit > 2000)
        .OrderByDescending(c => c.AccountOpenedDate)
        .Select(c => new { CustomerName = c.CustomerName, CreditLimit = c.CreditLimit });
}
\end{lstlisting}

Převod vrátil 33 záznamů: 27 \emph{Supplied} (jeden za každý fakt z~katalogu, např.\ \uv{Length 100 supplied by the database catalog from 'Sales.Customers'}), 4 \emph{Incompleteness} a 2 \emph{Convention}. Záznamy stojí za přečtení, protože říkají, kde výstup tvrdí víc než vstup:

\begin{itemize}[leftmargin=*,itemsep=1pt]
\item \emph{Incompleteness} ke klíči: \uv{The key column 'CustomerID' of 'Sales.Customers' is filled by a default constraint; the catalog cannot name the mechanism, so the generation strategy stays unspecified (decision 064).} -- a navazující \emph{Convention}: EF~Core generuje celočíselný klíč sám, výstup proto strategii nevyslovuje.
\item \emph{Incompleteness} ke vztahům: katalog zná cizí klíč \texttt{BillToCustomerID} tabulky \texttt{Customers} na sebe samu a cizí klíč z~\texttt{CustomerTransactions}, ale entity k~nim nemají navigační vlastnost, vztah se negeneruje. Naopak kolekce \texttt{Transactions}, kterou zdroj deklaruje, dostala vztah z~katalogu (\emph{Supplied}: \uv{the source declares the collection, the schema the relation}).
\end{itemize}

Kniha tvrzení má 27 řádků, všechny s~původem \emph{catalog} a osudem \emph{carried}. Bez katalogu by týž vstup dal entitu \texttt{[Keyless]} s~konvenčními jmény tabulek a záznamy \emph{Convention} -- a při cíli NHibernate nebo Hibernate \emph{Failure}, protože bez klíče entitu nepřijmou.

\subsubsection*{Příklad 2 (střední): EF Core → Hibernate, čtyři entity a šest dotazů v~LINQ přes hranici ekosystémů}

Příklad \emph{order-book} nese čtyři anotované entity EF~Core (zákazník, objednávka, řádek se složeným klíčem, produkt s~přiřazeným klíčem) a šest metod s~dotazy v~LINQ: stránkované otevřené objednávky s~\texttt{Contains} nad polem, nejprodávanější produkty se seskupením a \texttt{HAVING}, spící zákazníci s~\texttt{!Any}, objednávky s~produkty přes \texttt{IN} nad poddotazem, produkty nad průměrnou cenou se skalárním poddotazem a sjednocení dvou projekcí. Výpis~\ref{lst:ex2-in} ukazuje jednu entitu a dva dotazy.

\begin{lstlisting}[language={[Sharp]C},caption={Příklad 2, vstup (výběr): \texttt{OrderLine.cs}, \texttt{OpenOrders.cs}, \texttt{BestSellers.cs}.},label={lst:ex2-in}]
[Table("OrderLines", Schema = "Ordering")]
[PrimaryKey(nameof(SalesOrderId), nameof(LineNumber))]
public class OrderLine
{
    public int SalesOrderId { get; set; }
    public short LineNumber { get; set; }
    [ForeignKey(nameof(SalesOrderId))] public required SalesOrder Order { get; set; }
    public int ProductId { get; set; }
    [ForeignKey(nameof(ProductId))] public required Product Product { get; set; }
    public int Quantity { get; set; }
    [Precision(18, 2)] public decimal UnitPrice { get; set; }
    [Precision(5, 2)] public decimal DiscountPercent { get; set; }
}

public List<SalesOrder> OpenOrders(int customerId, DateTime placedAfter, int skip, int take)
{
    return ctx.SalesOrders
        .Where(o => o.CustomerId == customerId && o.OrderedAt >= placedAfter
                 && new[] { "New", "Paid", "Packed" }.Contains(o.Status))
        .OrderByDescending(o => o.OrderedAt).ThenBy(o => o.OrderNumber)
        .Skip(skip).Take(take).ToList();
}

public void BestSellers(int minimumQuantity)
{
    var rows = ctx.OrderLines
        .Where(l => l.DiscountPercent < 50)
        .GroupBy(l => l.ProductId)
        .Where(g => g.Sum(l => l.Quantity) >= minimumQuantity)
        .OrderBy(g => g.Key)
        .Select(g => new { ProductId = g.Key, Quantity = g.Sum(l => l.Quantity),
                           Lines = g.Count(), HighestPrice = g.Max(l => l.UnitPrice) })
        .ToList();
}
\end{lstlisting}

Výstup (výpis~\ref{lst:ex2-out}) je šestnáct artefaktů: čtyři entity Javy s~anotacemi \texttt{jakarta.persistence} a ke každému dotazu metoda nad \texttt{EntityManager} i holé JPQL. Složený klíč vyšel jako \texttt{@IdClass} s~vnořenou třídou \texttt{OrderLineId} s~\texttt{equals} a \texttt{hashCode}; cizí klíč, který je zároveň částí klíče, jako \texttt{@JoinColumn(insertable = false, updatable = false)} vedle klíčové části. Typy přešly slovníkem: \texttt{decimal} na \texttt{BigDecimal}, \texttt{DateOnly} na \texttt{LocalDate}, \texttt{[Precision(3)] DateTime} na \texttt{LocalDateTime} s~\texttt{secondPrecision = 3}, \texttt{byte[]} s~\texttt{[Timestamp]} na \texttt{@Version byte[]} (Hibernate 7.4.5 ji přijímá; EclipseLink by ji nepřijal, viz~\ref{sec:chyby-vypocet}).

\begin{lstlisting}[language=Java,caption={Příklad 2, výstup (výběr): \texttt{OrderLine.java} bez getterů a setterů, metody dotazů.},label={lst:ex2-out}]
@Entity
@Table(name = "OrderLines", schema = "Ordering")
@IdClass(OrderLine.OrderLineId.class)
public class OrderLine {
    @Id @Column(name = "SalesOrderId") private Integer SalesOrderId;
    @Id @Column(name = "LineNumber") private Short LineNumber;
    @Column(name = "ProductId", nullable = false, insertable = false, updatable = false) private int ProductId;
    @Column(name = "Quantity", nullable = false) private int Quantity;
    @Column(name = "UnitPrice", precision = 18, scale = 2, nullable = false) private BigDecimal UnitPrice;
    @Column(name = "DiscountPercent", precision = 5, scale = 2, nullable = false) private BigDecimal DiscountPercent;
    @ManyToOne(optional = false)
    @JoinColumn(name = "SalesOrderId", referencedColumnName = "SalesOrderId", nullable = false, insertable = false, updatable = false)
    private SalesOrder Order;
    @ManyToOne(optional = false)
    @JoinColumn(name = "ProductId", referencedColumnName = "ProductId", nullable = false)
    private Product Product;

    public static class OrderLineId implements Serializable {
        private Integer SalesOrderId;
        private Short LineNumber;
        /* konstruktory, equals, hashCode */
    }
}

public static TypedQuery<SalesOrder> query(EntityManager em, int customerId, LocalDateTime placedAfter, int skip, int take) {
    return em.createQuery("""
        select o
        from SalesOrder o
        where o.CustomerId = :customerId and o.OrderedAt >= :placedAfter and o.Status in ('New', 'Paid', 'Packed')
        order by o.OrderedAt desc, o.OrderNumber asc
        """, SalesOrder.class)
        .setParameter("customerId", customerId)
        .setParameter("placedAfter", placedAfter)
        .setFirstResult(skip)
        .setMaxResults(take);
}

public static Query query(EntityManager em, int minimumQuantity) {
    return em.createQuery("""
        select l.ProductId as ProductId, sum(l.Quantity) as Quantity, count(l) as Lines, max(l.UnitPrice) as HighestPrice
        from OrderLine l
        where l.DiscountPercent < 50
        group by l.ProductId
        having sum(l.Quantity) >= :minimumQuantity
        order by l.ProductId asc
        """)
        .setParameter("minimumQuantity", minimumQuantity);
}
\end{lstlisting}

Převod vrátil devět záznamů. Dva \emph{Loss}: neunikátní index nad \texttt{OrderedAt} (výkonnostní artefakt, mezireprezentace ho nenese, rozh.~055) a inicializátor \texttt{[]} u~\texttt{RowVersion}, který Java nepíše stejně. Čtyři \emph{Incompleteness} \uv{No table matching the entity was found in the catalog} -- instance měla katalog nad WideWorldImporters, které schéma \texttt{Ordering} nemá; to není chyba, jen zpráva, že doplnění nemělo co dodat. Tři \emph{Convention}: dvakrát \uv{the source leaves the mechanism to the dialect (Auto) \ldots the artifact writes IDENTITY, because AUTO would mean the implementation's own default -- a sequence or a counter table the database does not have (decision 187)} a jméno klíčové třídy \texttt{OrderLineId}. Jiný dotaz téhož příkladu ukazuje úpravu z~rozh.~155: skalární poddotaz \texttt{Average} vyšel jako \path{cast(avg(x.ListPrice) as BigDecimal(38,6))}, aby JPQL vrátilo tentýž typ jako LINQ.

\subsubsection*{Příklad 3 (složitý): EclipseLink → EF Core, jeden dlouhý dotaz se sedmi joiny a čtrnácti poddotazy}

Příklad \emph{marketplace-long-query} nese deset entit JPA tržiště a třídu \texttt{SalesDashboard}, jejíž metoda \texttt{regionalChampions} předává \texttt{EntityManager} dotaz v~JPQL o~sedmi joinech (jeden levý nad dvousloupcovou podmínkou), šestnácti agregátech v~projekci včetně \texttt{count(distinct)}, \texttt{case} a \texttt{coalesce}, skalárním poddotazu v~projekci, čtyřech úrovních vnoření v~podmínce (\texttt{in}, \texttt{exists}, \texttt{not exists}, \texttt{> any}, skalární \texttt{avg}), \texttt{having} se třemi podmínkami včetně poddotazu a stránkování (výpis~\ref{lst:ex3-in}, zkráceno na tvar).

\begin{lstlisting}[language=Java,caption={Příklad 3, vstup (výňatek): \texttt{SalesDashboard.java}.},label={lst:ex3-in}]
return em.createQuery("""
        select r.Name as Region, c.CustomerId as CustomerId, c.DisplayName as Customer,
               count(distinct o.OrderId) as Orders, count(ol) as Lines, sum(ol.Quantity) as Units,
               sum(ol.Quantity * ol.UnitPrice) as GrossRevenue,
               sum(case when ol.Discount is null then 0 else ol.Quantity * ol.UnitPrice * ol.Discount end) as DiscountGiven,
               coalesce(sum(rf.Amount), 0) as Refunded, avg(ol.UnitPrice) as AverageUnitPrice,
               ... ,
               (select count(rv3) from Review rv3 where rv3.CustomerId = c.CustomerId) as ReviewsWritten
        from OrderLine ol
        join SalesOrder o on o.OrderId = ol.OrderId
        join Customer c on c.CustomerId = o.CustomerId
        join Region r on r.RegionId = c.RegionId
        join Product p on p.ProductId = ol.ProductId
        join Category cat on cat.CategoryId = p.CategoryId
        join Supplier s on s.SupplierId = p.SupplierId
        left join Refund rf on rf.OrderId = ol.OrderId and rf.LineNumber = ol.LineNumber
        where o.Status in ('SHIPPED', 'DELIVERED')
          and o.OrderedAt between :fromDate and :toDate
          and lower(c.Email) not like '%!_test@%' escape '!'
          and cat.CategoryId in (select leaf.CategoryId from Category leaf
                where leaf.ParentCategoryId in (select branch.CategoryId from Category branch
                      where branch.ParentCategoryId is null and exists (select p2.ProductId from Product p2
                            where p2.CategoryId = leaf.CategoryId and p2.ListPrice > (select avg(p3.ListPrice) from Product p3
                                  where p3.SupplierId in (select s2.SupplierId from Supplier s2
                                        where s2.RegionId = c.RegionId and s2.Rating >= 4)))))
          and not exists (select rf2.RefundId from Refund rf2 join SalesOrder o2 on o2.OrderId = rf2.OrderId
                where o2.CustomerId = c.CustomerId and rf2.Amount > (select max(ol2.UnitPrice) from OrderLine ol2 where ol2.OrderId = o2.OrderId))
          and ol.UnitPrice > any (select p4.ListPrice from Product p4 where p4.CategoryId = p.CategoryId and p4.Discontinued = true)
        group by r.Name, c.CustomerId, c.DisplayName
        having sum(ol.Quantity * ol.UnitPrice) > :minRevenue and count(distinct o.OrderId) >= :minOrders
           and avg(ol.UnitPrice) >= (select avg(p5.ListPrice) from Product p5 where ...)
        order by GrossRevenue desc, r.Name asc, c.CustomerId asc
        """, Tuple.class)
    .setParameter("fromDate", fromDate) ... .setFirstResult(skip).setMaxResults(take).getResultList();
\end{lstlisting}

Převod dal dvanáct artefaktů: deset entit EF~Core, jedenáctou entitu \texttt{ProductTag}, kterou nástroj syntetizoval z~\texttt{@ManyToMany} mezi produktem a značkou (rozh.~005: spojovací tabulka jako explicitní entita s~\texttt{[PrimaryKey(nameof(ProductId),} \texttt{nameof(TagId))]} a dvěma \texttt{[ForeignKey]}), a jeden dotaz v~LINQ (výpis~\ref{lst:ex3-out}, zkráceno). Třída \texttt{SalesDashboard} do výstupu nejde: záznam \emph{Convention} říká, že \uv{hands a query over to the framework in its member 'regionalChampions', so it is read as the code around its queries and not as an entity}.

\begin{lstlisting}[language={[Sharp]C},caption={Příklad 3, výstup (výňatek): metoda dotazu v~LINQ pro EF Core.},label={lst:ex3-out}]
public static IQueryable Query(DbContext ctx, DateTime fromDate, DateTime toDate, short minRating,
                               decimal minRevenue, long minOrders, int skip, int take)
{
    return ctx.Set<OrderLine>()
        .Join(ctx.Set<SalesOrder>(), ol => ol.OrderId, o => o.OrderId, (ol, o) => new { ol, o })
        .Join(ctx.Set<Customer>(), t => t.o.CustomerId, c => c.CustomerId, (t, c) => new { t.ol, t.o, c })
        ...
        .LeftJoin(ctx.Set<Refund>(), t => new { t.ol.OrderId, t.ol.LineNumber },
                  rf => new { rf.OrderId, rf.LineNumber }, (t, rf) => new { t.ol, t.o, t.c, t.r, t.p, t.cat, t.s, rf })
        .Where(t => new[] { "SHIPPED", "DELIVERED" }.Contains(t.o.Status)
            && t.o.OrderedAt >= fromDate && t.o.OrderedAt <= toDate
            && !(t.c.Email.ToLower().Contains("_test@"))
            && ctx.Set<Category>().Where(leaf => ctx.Set<Category>().Where(branch => branch.ParentCategoryId == null
                 && ctx.Set<Product>().Where(p2 => p2.CategoryId == leaf.CategoryId && p2.ListPrice > ctx.Set<Product>()
                      .Where(p3 => ctx.Set<Supplier>().Where(s2 => s2.RegionId == t.c.RegionId && s2.Rating >= 4)
                           .Select(s2 => s2.SupplierId).Contains(p3.SupplierId)).Average(p3 => p3.ListPrice)).Any())
                 .Select(branch => branch.CategoryId).Contains(leaf.ParentCategoryId.Value))
                 .Select(leaf => leaf.CategoryId).Contains(t.cat.CategoryId)
            && ...
            && ctx.Set<Product>().Where(p4 => p4.CategoryId == t.p.CategoryId && p4.Discontinued == true)
                 .Select(p4 => p4.ListPrice).Any(v => t.ol.UnitPrice > v))
        .GroupBy(t => new { t.r.Name, t.c.CustomerId, t.c.DisplayName })
        .Where(g => g.Sum(t => t.ol.Quantity * t.ol.UnitPrice) > minRevenue
            && g.Select(t => t.o.OrderId).Distinct().Count() >= minOrders && ...)
        .Select(g => new { Region = g.Key.Name, CustomerId = g.Key.CustomerId, Customer = g.Key.DisplayName,
            Orders = g.Select(t => t.o.OrderId).Distinct().Count(), Lines = g.Count(),
            GrossRevenue = g.Sum(t => t.ol.Quantity * t.ol.UnitPrice),
            DiscountGiven = (g.Count(t => (t.ol.Discount == null ? 0 : t.ol.Quantity * t.ol.UnitPrice * t.ol.Discount) != null) > 0
                ? (decimal?)g.Sum(t => (t.ol.Discount == null ? 0 : t.ol.Quantity * t.ol.UnitPrice * t.ol.Discount)) : null),
            Refunded = g.Sum(t => t.rf.Amount), ...,
            ReviewsWritten = ctx.Set<Review>().Where(rv3 => rv3.CustomerId == g.Key.CustomerId).Count() })
        .OrderByDescending(p => p.GrossRevenue).ThenBy(p => p.Region).ThenBy(p => p.CustomerId)
        .Skip(skip).Take(take);
}
\end{lstlisting}

Na tomto příkladu je vidět, co znamená \emph{věrný} překlad a kde stojí jeho cena. Spojení je řetěz \texttt{Join} s~anonymními typy, \texttt{between} vyšlo jako dvojice porovnání (záznam \emph{Convention} s~pravidlem Q14), \texttt{like \ldots escape} jako \texttt{Contains} po odúnikování vzoru (rozh.~051), \texttt{> any} jako \texttt{Any(v => \ldots)} (rozh.~119), \texttt{count(distinct)} jako \texttt{Select(\ldots).Distinct().Count()}. Nejnápadnější je \texttt{DiscountGiven}: JPQL \texttt{sum} nad výrazem dává nad prázdnou skupinou \texttt{NULL}, zatímco \texttt{Sum} EF~Core nulu; builder proto píše podmíněný výraz, který vrací \texttt{null}, když skupina nemá žádnou hodnotu (rozh.~155) -- méně hezké, ale shodné řádky. Záznamy (29): deset \emph{Incompleteness} za entity, které katalog instance nezná; pět \emph{Loss} za inicializátory kolekcí; \emph{Loss} za \texttt{CHAR(2)} sloupec \texttt{CountryCode} (\uv{the source's framework hands the value \ldots without trailing spaces, and EFCore has no place in the entity artifact to trim it}, rozh.~169); \emph{Loss} za verzi, kterou Hibernate zvyšuje sám a EF~Core jen porovnává; a sedm záznamů ke strategiím klíčů, které EF~Core anotacemi nevyjádří (sekvence a \texttt{IDENTITY} jsou fluent-only, rozh.~187). Kniha tvrzení má 186 řádků, 178 nesených a 8 vypuštěných -- těch, k~nimž patří záznamy \emph{Loss}.

\subsection{Použití experimentů}\label{sec:exp-navod}

Experimenty jsou testy obou sad, které běží jen s~nastavenou proměnnou prostředí a zapisují soubory do adresáře běhu; bez proměnné přeskočí s~důvodem. Pracují nad třemi druhy vstupů (tabulka~\ref{tab:vstupy}), které si \emph{.NET sada sama vypíše} do \texttt{inputs/} adresáře běhu -- jeden podadresář na klíč a zdrojový framework s~jednotkami a \texttt{input.json}, u~klíčů s~katalogem i \texttt{schema.sql} s~DDL katalogu (rozh.~207). Tentýž výpis dostává jazykový model, takže se srovnává nad týmiž bajty.

\begin{table}[ht]
\centering\small
\begin{tabularx}{\textwidth}{L{2.4cm} Y L{3.6cm}}
\toprule
\textbf{Druh (klíč)} & \textbf{Co to je} & \textbf{Soudce 4.~stupně} \\
\midrule
\texttt{matrix:} & kategorie manifestu \path{Tests/Database/QueryShapes/categories.txt} nad doménou sedmi entit se složenými klíči (2--4 části) a vícesloupcovými cizími klíči, z~každého zdroje, který kategorii vysloví & kanonický výsledek v~repozitáři (\path{Tests/Database/Differential/matrix.txt}) \\
\texttt{cross:} & scénář překladu přes hranici ekosystémů ze šesti zdrojů, jen entity & -- (3.~stupeň) \\
\texttt{example:} & příklady výkladové stránky včetně šesti výňatků skutečných projektů (Keycloak, Cofoundry, Ambari, Chinook ×~2, JPetStore) & případová studie běží nad prázdnými tabulkami svého schématu (\emph{spustitelný}, bez ekvivalence); ostatní 3.~stupeň \\
\texttt{ldbc:} & 41 čtecích dotazů LDBC SNB Interactive v1 nad \texttt{LdbcSnb} & validační sada LDBC (zvlášť, soudce LDBC) \\
\texttt{fiben:} & 237 dotazů FIBEN s~otázkou v~přirozeném jazyce nad \texttt{Fiben} & soudce z~prvního doloženého běhu (regresní) \\
\texttt{sede:} & 1~714 dotazů SEDE s~titulkem a parametry nad \texttt{Sede} & soudce z~prvního dokončeného běhu (zatím není) \\
\bottomrule
\end{tabularx}
\caption{Vstupy experimentů.}
\label{tab:vstupy}
\end{table}

\paragraph{Experimentální běh (T2 × T3).} Tři kroky v~tomto pořadí, na hostiteli s~oběma toolchainy a testovací databází, nebo v~kontejnerech jako \texttt{docker compose --profile test run} s~touž proměnnou a adresářem jako svazkem:

\begin{lstlisting}[language=bash,caption={Experimentální běh.}]
export ORMCONVERTOR_EXPERIMENT_RESULTS=<adresář běhu>
dotnet test Tests/Tests.csproj --configuration Release --filter "FullyQualifiedName~Tests.Experiment.ExperimentRunTest"
(cd JavaTests && mvn -B test -Dtest=ExperimentRunTest)     # ORMCONVERTOR_API_URL instance nad testovací databází
dotnet test Tests/Tests.csproj --configuration Release --filter "FullyQualifiedName~ExperimentSummaryTest"
\end{lstlisting}

.NET půlka vypíše vstupy, přeloží každý do šesti cílů, tři .NET cíle prožene stupni (přečteno, zkompilováno, přijato frameworkem, běží, shoda řádků) do \texttt{tool-rows.tsv} a změří úplnost přenesených mapovacích faktů zpětným čtením artefaktu parserem cíle do \texttt{mapping-rows.tsv} (rozh.~173, 188); javovská půlka soudí tři javovské cíle do \texttt{java-tool-rows.tsv} (\texttt{javac}, továrna Hibernatu, parser EclipseLinku, vazba MyBatisu, běh přes JDBC); souhrn složí \texttt{summary.md} s~podíly po dvojicích ekosystémů, směrech, kategoriích a druzích vstupu. Volby: \path{ORMCONVERTOR_EXPERIMENT_KINDS} omezí druhy (\texttt{matrix,cross,example,ldbc,fiben}), klíče \texttt{fiben:} potřebují \path{ConnectionStrings:FibenDatabase} a pro Javu instanci nad \texttt{Fiben}, případové studie databázi \texttt{ORMConvertorCaseStudies} a pro Javu instanci nad ní (\path{ORMCONVERTOR_CASESTUDY_API_URL}). Žádný řádek běh neshodí: odmítnutí je řádek \emph{refused}, červený stupeň je řádek; shodí ho jen vada cesty (klíč bez měřidla, nenakonfigurovaná databáze druhu).

\paragraph{Srovnání s~jazykovým modelem.} Protokol (rozh.~175, 189, 202--207) staví na témže výpisu a témže soudci; model je \emph{producent cizího artefaktu} (rozh.~172). Jeden adresář běhu, \path{ORMCONVERTOR_LLM_RUN}, sedm kroků:

\begin{enumerate}[leftmargin=*,itemsep=1pt]
\item \textbf{Buňky.} \texttt{cells.tsv} s~hlavičkou \texttt{answers}, \texttt{source}, \texttt{target}, \texttt{variant} (a volitelně \texttt{attempts}, výchozí 3, a \texttt{set}); jeden řádek na buňku, např.\ \path{matrix:Filtering  Dapper  Hibernate  zero-shot  basic}. Varianty se šablonou: \texttt{zero-shot}, \texttt{few-shot} (tři ručně přeložené příklady ve směru buňky) a \texttt{schema} (zero-shot s~DDL katalogu).
\item \textbf{Prompty.} \path{WebAttemptTest.ThePromptsAreWrittenWhereAsked} vypíše vstupy, \path{attempts/<id>/prompt.md} na pokus a tabulku \texttt{attempts.tsv}. Prompt opakuje modelu pravidla: bez webu, bez nástrojů, odpověď jako bloky \texttt{FILE:} s~celými soubory a závěrečný \texttt{REPORT:}.
\item \textbf{Pokusy} jedním ze tří kanálů. \emph{Webové rozhraní}: obsluha ručně, nová konverzace na pokus, prompt vložený celý, odpověď zkopírovaná do \texttt{reply.md}, časy a tokeny do tabulky jen tam, kde je rozhraní ukazuje. \emph{API}: \path{ApiAttemptTest.TheAttemptsAreMadeWhereAsked} nad koncovým bodem kompatibilním s~OpenAI podle \texttt{api.json} (adresa, model, snímek, teplota, strop tokenů, ceny; klíč výrobce jen jménem proměnné prostředí). \emph{Claude Code}: \path{CodeAttemptTest.TheAttemptsAreMadeWhereAsked} volá \texttt{claude -p} bez nástrojů na předplatném podle \texttt{code.json}, s~proxy, která každý požadavek zachytí a pokus propustí, jen když nese prompt a změřený kontext (rozh.~205). Pokusy jdou po jednom, tabulka se píše po každém, přerušený běh pokračuje, kde skončil.
\item \textbf{Import.} \path{WebAttemptTest.TheRepliesAreImportedWhereAsked} udělá z~každé odpovědi adresář artefaktu s~\texttt{artifact.json}; odpověď bez souborů nebo se stopou hledání (jméno repozitáře, odkaz, deklarovaný nástroj) je odmítnuta a pokus se dělá znovu. Odpověď \texttt{zero-shot} a \texttt{schema} se soudí, jak přišla; u~ostatních variant import opraví obálku -- chybějící \texttt{import}/\texttt{using}, třída okolo metody, jeden veřejný typ na soubor -- a zapíše, co opravil (rozh.~204).
\item \textbf{Nástroj nad týmiž vstupy.} Experimentální běh s~\path{ORMCONVERTOR_EXPERIMENT_RESULTS=<run>}.
\item \textbf{Soud.} \path{./foreign-run.sh <run>/artifacts <run> [s na artefakt]} -- proces na artefakt s~mezí času; nebo \texttt{ForeignArtifactRunTest} té sady, která cíl vlastní, pro kontrolu.
\item \textbf{Souhrn.} \texttt{ExperimentSummaryTest} přidá do \texttt{summary.md} oddíl \emph{Foreign artifacts beside the tool over the same cells}: po producentu, modelu, kanálu a variantě podíl nástroje, podíl po pokusech a best-of-n nad týmiž buňkami, podíl opravených obálek, a buňky jednu po druhé.
\end{enumerate}

Kroky 2 a 4 nepotřebují databázi ani instanci. Co se u~každého běhu zapisuje do záznamu (kanál, šablony, otisk promptu každé buňky, hranice trénovacích dat, odmítnuté odpovědi, součty času, tokenů a ceny, u~lokálního modelu otisk souboru vah a digest obrazu serveru), říká \texttt{docs/runs/README.md}; nasazení lokálního modelu a podobu výsledku popisuje kap.~\ref{sec:llm} a reprodukční návod.

\paragraph{Ostatní běhy.} Běh FIBEN a SEDE (\texttt{FibenRunTest}, \texttt{SedeRunTest}; proměnné s~prefixem \texttt{FIBEN}/\texttt{SEDE}: připojení, \texttt{\_QUERIES} pro výběr dotazů jako kontrolu, \texttt{\_RESULTS} pro adresář, \texttt{\_MEASURE=0} bez měření), soudce LDBC (\texttt{LdbcValidationTest}; \path{ORMCONVERTOR_LDBC_VALIDATION_ROWS} pro prefix sady, prázdné pro celou), běh mezí (\texttt{LimitSearchTest}; \path{ORMCONVERTOR_SERIES_AXES} pro výběr os, \path{ORMCONVERTOR_SERIES_CATALOG=none} pro běh nasucho) a běh celého vstupu (\texttt{WholeInputTest}) mají stejnou stavbu: proměnná zapne případ, adresář dostane soubory, kontrola z~hosta se nezapisuje, záznam vzniká jen z~běhu v~zafixovaném prostředí.
% ===== 04-denik.local.tex =====
\section{Z~deníku vývoje}\label{sec:denik}

Deník projektu je repozitář sám: každá netriviální volba je číslované rozhodnutí s~kontextem, zvažovanými variantami, volbou a důsledky (\texttt{docs/decisions/}, 207 souborů ve vydání 2.3.0; číslo 038 je vynechané), každý běh sad a každé měření má datovaný neměnný záznam (\texttt{docs/runs/}), stav repozitáře k~datu zachycuje osm auditů (\texttt{docs/audits/}) a historii nesou commity a anotace značek vydání (\texttt{git tag -n99}). Tato kapitola z~toho vybírá: jak práce probíhala, co se při implementaci ukázalo jako zajímavé, co se nepovedlo napoprvé a kde se výsledek liší od specifikace.

\subsection{Jak práce probíhala}

Fork vznikl 15.~července 2026 nad prototypem o~zhruba 600 testech; stav při převzetí je zmražený v~\texttt{docs/baseline.md}. Od té doby přibylo 399 commitů (asi 400~tisíc vložených a 84~tisíc smazaných řádků včetně generovaných katalogů a zámků balíčků), sedm vydání (tabulka~\ref{tab:vydani}) a kód nástroje mimo testy má asi 95~tisíc řádků C\#. Sady měly v~běhu za vydání 2.3.0 (2.~11.~2026) 13~947 testů .NET a 4~910 testů Javy, obě zelené v~zafixovaném prostředí.

\begin{table}[ht]
\centering\small
\begin{tabularx}{\textwidth}{L{1.4cm} L{2.1cm} Y}
\toprule
\textbf{Vydání} & \textbf{Datum} & \textbf{Co přineslo} \\
\midrule
1.0 & 21.~8.~2026 & .NET část dotažená: složené klíče, vztahy, doplnění z~katalogu, diagnostika, dotazová matice tří frameworků, čtyři stupně ověření; dokumentační systém \\
1.1.0, 1.2.0 & 23. a 26.~8. & vícesouborový vstup, REST kontrakt hlídaný testem, vydání jako anotovaná značka; stav, který popsala specifikace \\
2.0.0 & 22.~9. & javovský ekosystém: Hibernate, EclipseLink a MyBatis, javovská sada, diferenční ověření, překlad přes hranici oběma směry \\
2.1.0 & 3.~10. & překlad dotazů úplný v~rozsahu katalogu: výrazy, parametry, mezivýsledky, rekurze, nativní SQL jako úniková cesta, katalog LDBC \\
2.2.0 & 12.~10. & katalog LDBC přeložený celý a souzený validační sadou; kniha tvrzení ke každému mapovacímu faktu; archiv ZIP v~prohlížeči \\
2.3.0 & 5.~11. & katalogy LDBC a FIBEN souzené celé nad stromem vydání, experimentální běh s~tabulkou T2 × T3, hledání mezí, srovnání s~jazykovými modely, případové studie, připnuté prostředí \\
\bottomrule
\end{tabularx}
\caption{Vydání projektu.}
\label{tab:vydani}
\end{table}

Rytmus práce určovaly tři věci. \textbf{Rozhodnutí před implementací}: změna, která něco volí, začíná souborem v~\texttt{decisions/} s~jedinou doporučenou variantou a dlouhým zdůvodněním proti požadavkům a proti článku; teprve pak kód. \textbf{Povinný uzavírací krok}: každá změna chování aktualizuje \texttt{architecture.md}, škrtá položku z~\texttt{open-items.md} a kde mění nárok, i~\texttt{traceability.md} a \texttt{subset.md} -- proto dokumentace nezastarala, ale také proto bylo 2.~10.~2026 nutné živé dokumenty seškrtat na čtvrtinu, protože se do nich dostala narace místo stavu. \textbf{Audity} v~přelomových okamžicích (po čtvrtém kroku, před a po každém vydání, celého repozitáře 18.~10., robustnosti 30.~10., před odevzdáním 1.~11.) s~nálezy číslovanými podle závažnosti; co z~nich plynulo, se rozdělilo na opravu, rozhodnutí a otevřenou položku.

Práce probíhala na dvou strojích s~různým vybavením: notebook s~Dockerem, kde jedině běží kontejnery, Advisor a všechna měření, a pracovní stolní počítač bez Dockeru s~lokálním SQL Serverem Express a JDK uvnitř IntelliJ, kde běží obě sady jako kontrola. Odtud pravidlo, že nic se neměří tam, kde to nelze měřit v~zafixovaném prostředí, a že zelený běh na druhém stroji je kontrola, nikdy verdikt.

\subsection{Implementační zajímavosti a rozhodnutí}\label{sec:zajimavosti}

\paragraph{Fakta o~frameworku se měří, než se o~nich rozhodne.} Pravidlo, které vzniklo z~chyby. Rozh.~148 zamítlo \uv{vždy levý join} pro navigaci přes referenci v~operandu LINQ na základě nezměřeného tvrzení o~provideru NHibernatu a vydávalo vnitřní join, který potichu zahazoval řádky (audit 18.~10., nález 3.2). Náhrada (rozh.~151) vznikla až po sondě: \texttt{NHibernateLinqProviderTest} čte plán dotazu NHibernatu nad namapovaným modelem a \texttt{EFCoreQueryAcceptance} bere SQL EF~Core přes \texttt{ToQueryString()}, na javovské straně jeden test běží s~logováním SQL frameworku. Od té doby každé rozhodnutí, které přijímá nebo zamítá variantu kvůli chování frameworku, cituje test, který to měří, a test zůstává vedle kódu, aby upgrade, který fakt změní, selhal tam. Příklady: EclipseLink 5.0.0 ořezává \texttt{CHAR(n)} (rozh.~169), odmítá \texttt{@Version byte[]} (179) a plný vnější join (180); SQL Server nerozplete korelovaný poddotaz s~\texttt{TOP(1)} a opakuje jeho joiny pro každý vnější řádek, proto EF~Core bere hodnotu jedné skupiny přes \texttt{Max()} (159).

\paragraph{Mezireprezentace nese to, co frameworky sdílejí, a nic víc.} Vztah leží na entitě, ne na vlastnosti, a entity na sebe ukazují jménem, ne referencí (rozh.~001); složený klíč je seřazený seznam částí se strategií generování za část a vykresluje se ploše (006), takže \texttt{@EmbeddedId}, \texttt{<composite-id class>} i \texttt{[PrimaryKey]} jsou jen formy téhož; N:M je explicitní spojovací entita (005), kterou nástroj syntetizuje i tam, kde ji zdroj nemá (příklad~3). Mezireprezentace nevaliduje -- připouští kombinace, které by databáze odmítla, a validace patří builderu cíle, který jediný ví, co framework unese. Podmínka dotazu je strom, ne plochý seznam konjunkcí; operand má šest tvarů (sloupec, literál, poddotaz, výčet hodnot, parametr, výraz; rozh.~024, 061, 074, 083, 107), a výraz stojí všude, kde stojí operand, včetně projekce a řazení.

\paragraph{Dotaz, který by vrátil jiné řádky, se nevydá.} Jediné pravidlo, které drží celou dotazovou větev pohromadě (rozh.~053, 065, 070): hranicí je \emph{množina řádků}. Co ji mění -- filtr, join, zdroj, seskupení, které parser nepřečte nebo cíl věrně nevyjádří -- odmítne artefakt; co mění jen to, co se o~řádku říká, je ztráta se záznamem. Z~pravidla plynou jednotlivosti: \texttt{IS NULL} je porovnávací operátor (002), \texttt{LIKE} se do LINQ překládá a ne přenáší (051), predikát LINQ nese trojhodnotovou logiku SQL oběma směry (156), \texttt{SUM} jazyka SQL a \texttt{Sum} EF~Core jsou dvě různé věci se dvěma zápisy (155), řádky se stejným klíčem řazení se porovnávají jako množina (170).

\paragraph{Javovská strana bez JVM.} Zásadní rozhodnutí cíle~2 (rozh.~076): javovské wrappery jsou v~C\# -- lexer celé lexikální gramatiky Javy a čtečka podmnožiny tříd, parser JPQL a vlastní zapisovač --, a JVM žije jen v~kontejneru testovací sady a na CI. Překladová cesta tak nemá druhý runtime a nasazení zůstává jeden proces; cenou je, že přijetí frameworkem a běh se ověřují jen testy. Hibernate a EclipseLink jsou dva tenké profily nad jednou vrstvou Jakarta Persistence (077, 080), které se liší výchozími hodnotami za týmiž anotacemi -- a právě ty výchozí hodnoty byly zdrojem překvapení: \texttt{AUTO} znamená u~jednoho sekvenci \texttt{<tabulka>\_SEQ}, u~druhého tabulku čítače, a žádná z~nich v~databázi konzumenta není, proto převod čte \texttt{AUTO} profilem zdroje a píše mechanismus nativní dialektu, \texttt{IDENTITY} (187). MyBatis nestojí na žádné vrstvě frameworku, ale na dvou vrstvách \emph{jazyka} -- Javě pro doménovou třídu a rozhraní, T-SQL pro dokument mapperu (084) --, a jeho \texttt{<id>} není klíč, ale identita řádku výsledku.

\paragraph{Javovská sada jako klient běžící instance.} Sada bere artefakty přes HTTP z~instance nástroje (rozh.~078), což dělá z~každého jejího testu zároveň test REST kontraktu a nutí compose i CI instanci spustit; na oplátku sada soudí přesně to, co dostane konzument. Diferenční ověření (089) spouští zdrojovou i přeloženou variantu dotazu a porovnává normalizované řádky s~jediným kanonickým výsledkem uloženým v~repozitáři, aby se dvě sady ve dvou jazycích shodly; negativní půlka mutuje artefakt a musí skončit rozdílem. Sada si sama počítá velikost a odmítne se sestavit pod kritériem F12 (087) -- počet testů tak není číslo v~dokumentu, které stárne, ale tvrzení, které drží sestavení.

\paragraph{Úniková cesta a druhý tvar.} Co jazyk cíle nevysloví, napíše cíl celým dotazem v~nativním SQL svého dialektu (rozh.~113). Rozhodnutí přišlo spolu s~katalogem LDBC, jehož dotazy potřebovaly rekurzi, seskupení podle výrazu, okenní funkce a agregaci do seznamu; bez únikové cesty by polovina katalogu končila odmítnutím, s~ní končí \emph{Fallback} se záznamem a měří se na čtvrtém stupni jako každý překlad. NHibernate píše každý dotaz dvakrát (118): závazným HQL a vedle něj LINQ nad \texttt{session.Query<T>()} sdíleným zapisovačem s~EF~Core; co provider 5.7.0 nevysloví, je \emph{Omitted}. Holé HQL čte vlastní sestupný parser (062), jediný úplný parser HQL totiž žije uvnitř NHibernatu.

\paragraph{LDBC a soudce, který není náš.} Katalog LDBC SNB Interactive v1 (rozh.~110) dal projektu referenční doménu s~nezávislou pravdou: validační sada LDBC se přehrává v~pořadí -- vkládání loaderem, každé čtení generovaným artefaktem každého frameworku -- a čtení se srovnává s~očekávanými řádky (117). Dotazy na cesty (IC~13, IC~14) nemají v~žádném ze šesti jazyků rekurzi, kterou by všechny unesly stejně; píší se proto jako prohledání do šířky rozepsané do řetězu nerekurzivních definic s~hloubkou, která je změřenou vlastností dat (140, 141) -- text dotazu nese číslo, které platí pro SF~1 a pro jiná data by se měřilo znovu. Soudce čte úsek mezi dvěma vkládáními ve všech frameworcích najednou, každý ve vlastním vlákně (166), což zkrátilo celou sadu z~několika dní na asi třicet hodin a dalo serveru druhý režim (164).

\paragraph{Měřicí prostředí a politika iterací.} Čísla jsou obhajitelná jen nad zafixovaným prostředím (rozh.~145): SQL Server s~pevnou pamětí a bez paralelismu, databáze a klienti na oddělených procesorech, virtuální stroj se čtyřmi procesory a 12~GB bez swapu (193), turbo hostitele vypnuté stavem procesoru 99~\%; krátký dotaz se neprodlužuje, opakuje se do meze a vykazuje se medián, buňka se zahřívá jedním studeným spuštěním (165). Soudce FIBEN (149) je jiného druhu než LDBC: vznikl z~prvního doloženého běhu a každý další běh se s~ním srovnává -- regrese mezi verzemi, ne správnost čtení otázky; nový soudce je rozhodnutí s~datem a důvodem (157, 163), nikdy oprava v~porovnání.

\paragraph{Robustnost.} Audit robustnosti (30.~10.) a běh mezí (182) odkryly, kde nástroj končí: rekurzivní sestup nad hluboce zanořeným vstupem přetékal zásobník, což .NET nechytí a proces zemře. Strop hloubky zanoření se počítá nad tokeny ještě před sestupem (092) a týmž číslem i nad modelem po přečtení -- řetězy bez závorek, \texttt{CASE} až po \texttt{END} -- a čtečka, které dochází zásobník, jednotku odmítne (191); text jednoho dotazu má strop délky, který drží i operand LINQ, rozvinutí \texttt{<include>} a hotový artefakt (194); řetěz HQL nad 32 článků jde nativním SQL, protože tam končí vlastní parser NHibernatu (192). Běh mezí pak měří, kde se zastaví překladač a kde cesta ověření: Roslyn kompiluje vnořené lambdy exponenciálně v~hloubce, parser HQL NHibernatu přetéká na 1024 poddotazech, generátor SQL EF~Core na 4096 výrazech -- meze cílů, ne nástroje.

\paragraph{Cizí artefakt a protokol srovnání.} Aby se dal jazykový model soudit týmž metrem jako nástroj, vstupuje jeho výstup jako \emph{cizí artefakt}: adresář souborů s~manifestem, který říká, na který vstup odpovídá, komu patří a s~jakými metadaty (rozh.~172); sada ho prožene stupni ověření ve vlastním procesu s~mezí času a zapíše řádek s~původem. Protokol srovnání (175) nad tím stojí: táž vypsaná sada vstupů, tři pokusy na buňku, varianty a ablace jako hodnoty manifestu, podíly po pokusech i best-of-n. Tři kanály pokusů (webové rozhraní, API, Claude Code) vznikly postupně s~tím, jak se ukazovalo, co lze měřit: ve webovém rozhraní měří čas obsluha a model o~sobě lže (Qwen se hlásil jako GPT-4 i Claude), přes API jsou čas a tokeny z~koncového bodu, přes Claude Code dostává model kontext rozhraní, a proto ovladač každý požadavek zachytí proxy a pokus propustí, jen když nese prompt a změřený kontext (205). Odpověď modelu je celý soubor a soudí se jako celek vedle holé metody nástroje; u~variant mimo \texttt{zero-shot} import opraví obálku -- jen to, co konzument dává holé metodě -- a zapíše, co opravil (204), aby se měřil obsah překladu, ne shoda s~předepsaným tvarem.

\paragraph{Prostředí připnuté na obsah.} Zámek balíčků u~každého projektu, základní obrazy digestem, akce CI commitem, posun navrhuje Dependabot (rozh.~186); verze nástroje stojí na jediném místě (034) a číslo vydání se rozhoduje jednou za vydání podle toho, co nástroj umí, ne podle rozbité plochy (098) -- s~vysloveným předpokladem, že nástroj nemá konzumenta mimo repozitář. Žádný generovaný zdroj ani text o~běhu nenese přihlašovací údaje: harness Advisoru bere řetězec za běhu a cizí text odchází jedním začerněním (160).

\subsection{Co se nepovedlo napoprvé}\label{sec:nepovedlo}

\begin{itemize}[leftmargin=*,itemsep=2pt]
\item \textbf{Vnitřní join místo levého} (rozh.~148, výše): jediná známá vada, která tiše měnila množinu řádků; našel ji audit, ne test, protože matice tehdy neměla kategorii s~chybějícím řádkem reference. Od té doby je v~matici.
\item \textbf{Soudce LDBC přerušený po 29 hodinách} (24.~10.): skript, který měl po skončení běhu vrátit databázi do měřicího režimu, se spustil, když \texttt{docker wait} ještě držel soubor přesměrování; \texttt{cmd} příkaz přeskočil, skript šel na \texttt{up -d mssql\_db} bez \texttt{judge.env}, databáze se znovu vytvořila a .NET soudce zemřel na řádku 128~816 ze 138~747 bez verdiktu. Pravidlo od té doby: skript se databáze dotkne až poté, co sám přečetl kontejner běhu jako \texttt{exited}.
\item \textbf{Zabité shelly pro nedostatek paměti} (22.~10., 5.~11.): dva SQL Servery a pět aplikačních kontejnerů najednou spotřebovaly 12~GB virtuálního stroje i 16~GB hostitele a hostitelský proces byl ukončen. Odtud pořadí \uv{zastavit systém, odstranit zbylé kontejnery, zkontrolovat volnou paměť} před každým během profilu \texttt{test} a rozh.~193.
\item \textbf{Javovská půlka běhu mezí padla haldou} (31.~10.): \texttt{OutOfMemoryError} jednoho bodu zabil vlákno HTTP klienta a každý další bod byl falešně \emph{crashed}. Sada od té doby bere 75~\% paměti kontejneru a na OOM se ukončí (\path{-XX:+ExitOnOutOfMemoryError}), místo aby běžela dál s~mrtvým vláknem; procházka pokračuje, kde skončila (195).
\item \textbf{První běh SEDE bez výsledku} (8.~11.): oba běhy držely kanonické řádky všech buněk v~paměti až do konce a soubory psaly až po posledním dotazu; FIBEN s~237 dotazy to unesl, SEDE s~1~714 dotazy, z~nichž řada vrací přes 100~tisíc řádků, ne -- .NET půlka padla po 8~h na poslední alokaci při skládání textu, javovská byla zastavena ve třetině s~haldou u~stropu. Náprava je v~kódu (soubory běhu se píší po dotazech a řádky se po porovnání pouštějí), soudce SEDE vznikne z~příštího dokončeného běhu.
\item \textbf{Výňatky, které se nezkompilovaly} (24.~10.): první kontrola experimentálního běhu nad příklady ukázala, že výňatky skutečných projektů sahají na typy mimo výňatek, výčty a mapy; vstupy byly upraveny a kontrola, která změnila vstupy experimentu, je proto výjimečně zapsaná jako běh. Výňatek Cofoundry (druhá aplikace T1, rozh.~184) i po úpravě končí v~prvním stupni: jeho dotazy sahají na typ \texttt{Template} mimo výňatek a člen zachyceného objektu se četl jako sloupec -- otevřená položka.
\item \textbf{Zacyklený pilotní model} (1.~11.): třetina odpovědí Qwen3.5-9B skončila stropem tokenů, protože se v~poli \texttt{assumptions} zprávy opakovala jedna věta až 130krát; každý takový pokus stál hodinu místo deseti minut. Je to výsledek měřený protokolem (teplota 0 bez penalizace opakování), ne nezdar kanálu, a záznam to tak nese.
\end{itemize}

\subsection{Odchylky od specifikace}\label{sec:odchylky}

\begin{itemize}[leftmargin=*,itemsep=2pt]
\item \textbf{Mezireprezentace se uživateli nezobrazuje.} Specifikace (F14, milník~4) slibovala zobrazení mezireprezentace; rozh.~154 to zavrhlo: uživateli by ukázala vnitřní slovník nástroje, zatímco kniha tvrzení a záznamy říkají totéž slovy vstupu a výstupu. Nárok F14 to říká výslovně.
\item \textbf{Obrazovka bez číslovaných kroků.} Makety specifikace měly pět očíslovaných kroků; dnešní obrazovka má oddíly bez čísel, volitelný příklad jako první oddíl a vstup vedle výstupu párovaný po tvrzeních (rozh.~135, 137). Scénář S7 zůstává pod pěti kroky.
\item \textbf{Advisor nebyl rozšířen} (F15, T7): podle rozh.~171 přijde až po experimentech, zbude-li čas; čas nezbyl. Měří dál jen Dapper a EF~Core a je vyňatý ze záruk, což stránky říkají samy. Javovská větev Advisoru je rozhodnutá (146), ne implementovaná.
\item \textbf{T1 v~užším smyslu}: nárokuje se vstup experimentu (výňatek aplikace splňuje čísla zadání, jak ho přečte převod), ne výsledek nad ním; případová studie je nad prázdnými tabulkami jen \emph{spustitelná}, bez měřítka ekvivalence.
\item \textbf{Víc testovacích sad, než specifikace jmenovala}: katalogy LDBC, FIBEN a SEDE nejsou ve specifikaci; vznikly, protože matice vlastních kategorií nestačila jako důkaz praktické použitelnosti a protože srovnání s~LLM potřebuje vstupy, které nástroj nepsal.
\item \textbf{LLM experimenty byly částečně provedeny}, ač nebyly podmínkou: protokol, tři kanály a čtyři záznamy srovnání (kap.~\ref{sec:llm}); plné zero-shot, few-shot, RAG a agentní varianty zůstávají diplomové práci, jak specifikace říká.
\item \textbf{Javovská sada má třetí běhové prostředí}: vedle kontejneru a CI běží i na stolním počítači z~JDK a Mavenu uvnitř IntelliJ -- jako kontrola, ne verdikt; věta specifikace \uv{JVM žije jen v~kontejneru a na CI} platí pro překladovou cestu a pro měření.
\end{itemize}

\subsection{Využití AI}\label{sec:ai}

Při vývoji jsme soustavně používali asistenta pro programování (Claude Code od společnosti Anthropic) jako pracovní nástroj pod stejným režimem, jaký platí pro každou jinou změnu v~repozitáři: netriviální volba začíná návrhem rozhodnutí s~jedinou doporučenou variantou a zdůvodněním proti požadavkům a článku, které řešitel přijme, upraví nebo zamítne; teprve pak vzniká kód, testy a aktualizace dokumentace. Asistent psal podle přijatých rozhodnutí kód wrapperů, testů a dokumentace, připravoval audity a záznamy běhů a navrhoval znění této zprávy z~dokumentů repozitáře. Řešitel nese volbu každého rozhodnutí, čte a schvaluje změny, spouští a hodnotí běhy sad a měření v~zafixovaném prostředí, commituje a vydává. Dvě pravidla z~toho vzešla přímo: fakta o~chování frameworku se měří testem proti připnuté verzi dřív, než se o~nich rozhodne (kap.~\ref{sec:zajimavosti}; rozh.~148 je případ, kdy nezměřené tvrzení prošlo), a nic, co nelze na daném stroji spustit, se nepovažuje za ověřené. Ve srovnání s~jazykovými modely (kap.~\ref{sec:llm}) je týž produkt zároveň předmětem měření; oba role jsou oddělené: kanál \texttt{code} volá model bez nástrojů, v~prázdné složce mimo repozitář a s~každým požadavkem zachyceným proxy, takže k~modelu nedojde nic z~kontextu vývoje (rozh.~205), a prompt repozitář ani nástroj nejmenuje.
% ===== 05-experimenty.local.tex =====
\section{Experimentální část}\label{sec:exp}

Specifikace rozdělila experimenty na požadavky projektu (T1--T3: případová studie, matice překladů, metriky korektnosti) a navazující práci (T4--T6: LLM baseline, RAG a agentní varianty, ablace; T7: Advisor). Projekt T1--T3 nárokuje v~užším smyslu -- vstup, matici a \emph{měřidlo}, ne čísla (rozh.~183, 184) -- a čísla níže jsou záznamy verze 2.3.0 v~\texttt{docs/runs/}, každý s~commitem, strojem a režimem serveru. T4--T6 nejsou nárokované; protokol a čtyři záznamy srovnání (kap.~\ref{sec:llm}) ukazují, že měřicí infrastruktura stojí a co na vzorcích změřila.

\subsection{Metodika}

\textbf{Stupně ověření} (rozh.~016, 027): (1)~artefakt je přečten parserem jazyka cíle, (2)~zkompilován (Roslyn, \texttt{javac}, validace XML), (3)~přijat frameworkem (továrna NHibernatu, model EF~Core, \texttt{SessionFactory} Hibernatu, parser EclipseLinku, vazba MyBatisu), (4)~spuštěn proti databázi a jeho řádky se shodují s~měřítkem. \textbf{Metriky T3} jsou podíly artefaktů na každém stupni a úplnost přenesených mapovacích faktů: podíl faktů reference, které parser \emph{téhož} frameworku přečte ze svého artefaktu zpět, s~hodnotou uznanou i tam, kde fakt vyslovuje jinými slovy (\texttt{Identity} za \texttt{Auto}, klíč za NOT NULL svých sloupců; rozh.~188). \textbf{Měřítka 4.~stupně} jsou tři podle druhu vstupu (tab.~\ref{tab:vstupy}): kanonický výsledek v~repozitáři u~matice, nezávislá validační sada u~LDBC, regresní soudce z~doloženého běhu u~FIBEN. \textbf{Prostředí} je zafixované (rozh.~145, 164, 193) a každý záznam říká, co platilo. \textbf{Populace} experimentálního běhu za 2.3.0: 624 vstupů výpisu, každý do šesti cílů -- matice 318 (59 kategorií z~každého zdroje, který ji vysloví), FIBEN 237, LDBC 41, příklady 22 (z~nich šest případových studií), scénáře \texttt{cross} 6; celkem 3~744 řádků nástroje a 3~744 řádků úplnosti mapování.

\subsection{Výsledky nástroje}

\paragraph{Matice T2 × T3.} Tabulka~\ref{tab:t2t3} je záznam experimentálního běhu z~5.~11.~2026 (strom \texttt{a105b3d}) po dvojicích ekosystémů: \emph{n} překlady, \emph{odmítnuto} bez artefaktu, \emph{nativní SQL} úniková cesta, stupně 1--3, \emph{n4} překlady se 4.~stupněm, \emph{běží} a \emph{shoda} z~nich, \emph{mapování} úplnost přenesených faktů. Každá z~59 kategorií matice je ve všech čtyřech dvojicích přijatá, běží a shoduje se s~kanonickým výsledkem; FIBEN je ve shodě se soudcem v~každém Kimově typu a každém cíli; LDBC končí třetím stupněm (soudí ho validační sada zvlášť) bez odmítnutí.

\begin{table}[ht]
\centering\small
\setlength{\tabcolsep}{4pt}
\begin{tabular}{l r r r r r r r r r r}
\toprule
\textbf{Dvojice} & \textbf{n} & \textbf{odm.} & \textbf{nat.\ SQL} & \textbf{přečteno} & \textbf{zkomp.} & \textbf{přijato} & \textbf{n4} & \textbf{běží} & \textbf{shoda} & \textbf{mapování} \\
\midrule
.NET → .NET & 1341 & 1 & 129 & 1340 & 1338 & 1337 & 1186 & 1183 & 1179 & 99 \% \\
.NET → Java & 1341 & 2 & 95 & 1339 & 1336 & 1335 & 1186 & 1183 & 1179 & 69 \% \\
Java → .NET & 531 & 4 & 42 & 527 & 527 & 527 & 492 & 489 & 486 & 99 \% \\
Java → Java & 531 & 4 & 29 & 527 & 527 & 527 & 492 & 490 & 486 & 75 \% \\
\bottomrule
\end{tabular}
\caption{Tabulka T2 × T3 po dvojicích ekosystémů, experimentální běh 5.~11.~2026.}
\label{tab:t2t3}
\end{table}

Jedenáct odmítnutí je podle rozhodnutí: zdroje Dapper a MyBatis bez katalogu do cílů, které klíč vyžadují, JPetStore do EF~Core (sloupce výsledku bez typu), Keycloak do EF~Core (kolekční parametr v~nativním SQL). Sedm červených řádků pod 4.~stupněm patří jednomu vstupu, výňatku Cofoundry (kap.~\ref{sec:nepovedlo}), a klíči, který je zároveň cizím klíčem, u~EclipseLinku (\texttt{marketplace-fluent}; vyřešeno po běhu rozh.~206). Nižší úplnost mapování javovských cílů (Hibernate 66~\%, EclipseLink 72~\%, MyBatis 80~\%) dávají fakta, která parser JPA ani MyBatisu ze svého artefaktu nečte zpět -- databázový typ, který JPA nechává javovému typu, a \texttt{<id>} MyBatisu --, jménem držená testem jako vyslovená mez míry; .NET cíle jsou na 99~\%, Dapper nemá jmenovatele (deskriptor nevyjadřuje žádný mapovací fakt).

\paragraph{Případové studie (T1).} Tabulka~\ref{tab:studie}: stupeň, kterého projekt studie dosáhl jako celek v~každém cíli nad prázdnými tabulkami svého schématu; Dapper a MyBatis končí třetím stupněm záměrně (SQL bez schématu). Obě aplikace nárokované T1 -- Keycloak (Hibernate, 22 entit, 47 pojmenovaných dotazů) a Cofoundry (EF~Core, 29 entit) -- splňují čísla zadání, jak je z~výňatku spočítá sestavení; co nad nimi nástroj dokáže, je míra, a u~Cofoundry je nízká.

\begin{table}[ht]
\centering\small
\begin{tabular}{l c c c c c c}
\toprule
\textbf{Studie (zdroj)} & \textbf{Dapper} & \textbf{EF Core} & \textbf{NHibernate} & \textbf{Hibernate} & \textbf{EclipseLink} & \textbf{MyBatis} \\
\midrule
Keycloak (Hibernate) & 3 & odmítnuto & 4 & 4 & 4 & 3 \\
Cofoundry (EF Core) & 1 & 1 & 2 & 1 & 1 & 1 \\
Ambari (EclipseLink) & 3 & 4 & 4 & 4 & 4 & 3 \\
Chinook (Dapper) & 3 & 4 & 4 & 4 & 4 & 3 \\
Chinook (EF Core) & 3 & 4 & 4 & 4 & 4 & 3 \\
JPetStore (MyBatis) & 3 & odmítnuto & odmítnuto & odmítnuto & odmítnuto & 3 \\
\bottomrule
\end{tabular}
\caption{Případové studie: dosažený stupeň projektu v~každém cíli (5.~11.~2026).}
\label{tab:studie}
\end{table}

\paragraph{Katalogy.} Běh FIBEN za vydání 2.3.0 (2.~11.) dal v~obou sadách každou buňku ve shodě se zdrojem i se soudcem, bez selhání; z~237 dotazů jde do EF~Core 81, do NHibernatu 90 a do EclipseLinku 109 nativním SQL, do Hibernatu 15. Soudce LDBC nad celou validační sadou SF~1 (138~747 řádků, 7~041 vkládání; 2.--5.~11.) dal každé čtení ve shodě ve všech šesti frameworcích; přehrání trvalo 29,6~h (.NET) a 27,2~h (Java) v~režimu soudce, průměrný select čtení 630--735~ms, dotaz na cestu IC~14 asi 2~s. Oba běhy také odmítly pět mutací IC~2 v~každém frameworku. První běh SEDE (8.~11.) skončil bez výsledku (kap.~\ref{sec:nepovedlo}); soudce SEDE dá příští dokončený běh.

\paragraph{Meze a výkon.} Běh mezí (29.--30.~10.) prošel 8~963 bodů .NET půlky v~patnácti osách a 1~077 řádků javovské: překladač došel ke stropu řady ve všech směrech kromě tří -- řetěz mezivýsledků ze zdroje EF~Core roste exponenciálně v~délce (32 definic přes 35~s a 8~GB), 512 joinů ze zdroje EF~Core a projekt mezi 1~600 a 6~400 entitami překročily mez S3 (30~s); kde se zastavila cesta ověření, byly to meze cílů (Roslyn nad vnořenými lambdami, parser HQL NHibernatu, generátor SQL EF~Core), ne nástroje. Běh celého vstupu (30.~10.) sloučil všech 607 tehdejších vstupů do jednoho projektu na framework a přeložil šest projektů do šesti cílů: každý z~36 překladů pod třetinou meze S3 i s~katalogem, bez katalogu pod desetinou. Měření S3 nad testovacím projektem o~100 entitách a 100 dotazech dává 0,2--0,7~s, z~toho čtení katalogu 210--330~ms; podíl správně přečtených metadat F4 nad oběma schématy testovací databáze je 462/462.

\subsection{Srovnání s~jazykovými modely}\label{sec:llm}

Protokol (kap.~\ref{sec:exp-navod}) byl proveden ve čtyřech zaznamenaných bězích, tabulka~\ref{tab:llm-behy}. Vzorek prvních tří je týchž třináct buněk: kategorie matice ve všech čtyřech směrech (nejjednodušší filtrace v~každém směru, řazení do cíle se dvěma tvary, join nad dvěma sloupci do SQL, \texttt{CASE} v~projekci do mapperu, kolekční parametr \texttt{<foreach>} do \texttt{IN @ids}, korelovaný \texttt{EXISTS}, agregace do seznamu, kterou jazyk cíle nevysloví), jeden scénář jen s~entitami a jeden malý příklad; čtvrtý běh přidal 43 nových buněk ve třech sadách -- \emph{základní} (kategorie, které vzorek neměl), \emph{velká} (případové studie jako celé projekty, dlouhý dotaz, FIBEN s~otázkou, LDBC) a \emph{záludná} (kategorie a příklady, na nichž chyboval sám nástroj a oprava po sobě nechala test s~měřeným faktem frameworku).

\begin{table}[ht]
\centering\small
\begin{tabularx}{\textwidth}{L{1.8cm} L{3.6cm} L{1.5cm} Y}
\toprule
\textbf{Záznam} & \textbf{Model} & \textbf{Kanál} & \textbf{Vzorek, varianty, nastavení} \\
\midrule
1.~11. & Qwen3.5-9B, Q4\_K\_M (otevřené váhy, lokálně v~llama.cpp na CPU) & \texttt{api} & 13 buněk, \texttt{zero-shot}, teplota 0, seed 1, strop 8~192 tokenů; 39 pokusů, 17,1~h strojového času, 0~USD \\
3.~11. & Claude Fable 5.1 & \texttt{code} & 13 buněk, \texttt{zero-shot} a \texttt{few-shot}; 78 pokusů, teplota nenastavitelná, úsilí \texttt{high}; předplatné, cena podle ceníku 19,61~USD \\
5.~11. & Claude Fable 5.1 & \texttt{api} & tytéž buňky a prompty jako 3.~11.; 78 pokusů, 61,5~min, 21,68~USD účtovaných \\
6.~11. & Claude Opus 5.5 & \texttt{code} & 43 buněk ve třech sadách, obě varianty; 258 pokusů, dvě pětihodinová okna předplatného, cena podle ceníku 111,53~USD \\
\bottomrule
\end{tabularx}
\caption{Zaznamenané běhy srovnání s~jazykovým modelem.}
\label{tab:llm-behy}
\end{table}

\paragraph{Nasazení pilotního modelu.} Pilotní model je nejsilnější model s~otevřenými váhami, který se vejde na stroj se 16~GB (rozh.~202): \path{Qwen3.5-9B-Q4_K_M.gguf} z~\path{unsloth/Qwen3.5-9B-GGUF}, 5,68~GB, otisk SHA-256 v~záznamu, v~serveru llama.cpp z~obrazu \path{ghcr.io/ggml-org/llama.cpp:server} připnutém digestem. Na čtyřech jádrech virtuálního stroje dává asi 13 tokenů/s na promptu a 2,7 tokenu/s při generování, takže jeden pokus nad buňkou matice (prompt asi 2~000 tokenů, odpověď 1~500) trvá deset minut a nic jiného z~profilu \texttt{test} při něm neběží:

\begin{lstlisting}[language=bash,caption={Nasazení lokálního modelu pro kanál \texttt{api}.}]
sha256sum <models>/Qwen3.5-9B-Q4_K_M.gguf
docker pull ghcr.io/ggml-org/llama.cpp:server
docker image inspect -f '{{index .RepoDigests 0}}' ghcr.io/ggml-org/llama.cpp:server
MSYS_NO_PATHCONV=1 docker run -d --name llm_server -p 8080:8080 -v "<models>:/models:ro" \
  ghcr.io/ggml-org/llama.cpp:server -m /models/Qwen3.5-9B-Q4_K_M.gguf --host 0.0.0.0 --port 8080 \
  -c 16384 -t 4 --load-mode none --temp 0 --seed 1 --reasoning off
curl -s http://localhost:8080/health          # {"status":"ok"} po načtení, asi minuta
\end{lstlisting}

\texttt{--reasoning off}, protože se karta modelu a vydavatel GGUF neshodují na výchozím stavu uvažování; \texttt{--load-mode none}, protože model mapovaný z~windowsovského svazku stránkuje pomalu. \texttt{api.json} běhu jmenuje adresu \path{http://localhost:8080/v1}, model, otisk jako \texttt{version}, teplotu 0, seed 1, strop 8~192, ceny 0 a \path{cache_prompt: false}, aby druhý pokus buňky nedostal prompt z~cache a jeho čas byl časem celého pokusu. Pro model výrobce se v~témže souboru mění jen adresa, model, strop (32~768, protože se do něj počítá skryté uvažování), ceny ceníku a jméno proměnné prostředí s~klíčem; silný model výrobce teplotu nepřijímá, a proto jeho tři pokusy na buňku nejsou bajtově shodné (druhé nastavení rozh.~175).

\paragraph{Výsledky.} Tabulka~\ref{tab:llm-vysledky} staví vedle sebe podíl po pokusech a best-of-n na nejvyšším měřítku, které buňka má (\emph{shoda} u~matice a FIBEN, \emph{přijato} u~buněk, které končí třetím stupněm); nástroj dosáhl nejvyššího měřítka ve všech 13 buňkách prvních tří běhů a ve 41 ze 43 čtvrtého (dvě jsou odmítnutí podle rozhodnutí).

\begin{table}[ht]
\centering\small
\setlength{\tabcolsep}{4pt}
\begin{tabular}{l l l r r r r r}
\toprule
\textbf{Model, kanál} & \textbf{Varianta} & \textbf{Buněk} & \textbf{přečteno} & \textbf{zkompilováno} & \textbf{přijato} & \textbf{shoda} & \textbf{best-of-n shoda} \\
\midrule
Qwen3.5-9B, api & zero-shot & 13 & 8 \% & 8 \% & 0 \% & 0 \% & 0 \% \\
Fable 5.1, code & zero-shot & 13 & 100 \% & 100 \% & 87 \% & 85 \% & 91 \% \\
Fable 5.1, code & few-shot & 13 & 100 \% & 100 \% & 87 \% & 85 \% & 91 \% \\
Fable 5.1, api & zero-shot & 13 & 100 \% & 100 \% & 87 \% & 85 \% & 91 \% \\
Fable 5.1, api & few-shot & 13 & 100 \% & 100 \% & 87 \% & 85 \% & 91 \% \\
Opus 5.5, code & zero-shot & 43 & 98 \% & 89 \% & 73 \% & 77 \% & 89 \% \\
Opus 5.5, code & few-shot & 43 & 98 \% & 93 \% & 79 \% & 80 \% & 85 \% \\
\bottomrule
\end{tabular}
\caption{Podíly pokusů jazykových modelů po stupních (\emph{shoda} nad buňkami s~měřítkem ekvivalence; best-of-n podíl buněk, v~nichž uspěl aspoň jeden pokus).}
\label{tab:llm-vysledky}
\end{table}

\paragraph{Co z~toho plyne.}
\begin{itemize}[leftmargin=*,itemsep=2pt]
\item \textbf{Lokální 9B model drží tvar odpovědi, ne obálku.} Všech 39 odpovědí Qwenu nese bloky šablony, ale dotaz píše jako celý soubor tam, kde prompt (šablona \texttt{web-zero-shot-2}) žádal holou metodu, a mapování NHibernatu bez tříd entit; proto 0 přijatých. Kontrola mimo záznam, která odpovědi opravila jen mechanicky (obálka, \texttt{[PrimaryKey]} místo několika \texttt{[Key]}, schéma do \texttt{[Table]}), dovedla 8 ze 13 buněk na nejvyšší stupeň; 5 zůstalo na chybách obsahu -- špatná projekce agregace, T-SQL vložené do HQL, vymyšlené parametry, zástupné symboly pojmenované hodnotou. Zjištění vedlo k~rozh.~204 (odpověď jako celý soubor) a varianta \texttt{zero-shot} se od té doby soudí, jak přišla, zatímco ostatní varianty dostanou obálku opravenou.
\item \textbf{Silný model přeloží obsah, ale ne fakta, která nezná.} U~Fable 5.1 projde každý ze 78 artefaktů čtením i kompilací a jediná příčina každého neúspěchu je atribut \texttt{assembly} v~mapování NHibernatu: deset z~dvanácti pokusů ho napíše odhadem a soudce, který jméno sestavení dodává sám (rozh.~028), mapování odmítne. Kanál \texttt{code} a \texttt{api} dávají týž obraz; liší se tím, co model ví o~sobě (z~kontextu Claude Code zná svou verzi, přes API ne) a účtováním. \texttt{few-shot} u~tohoto modelu verdikt nemění: oprav obálky bylo 25 z~39 a všechny jsou chybějící \texttt{import}/\texttt{using}.
\item \textbf{Nad složitějšími vstupy roste rozdíl.} Opus 5.5 nad \emph{základní} sadou chybuje jen atributem \texttt{assembly}; nad \emph{záludnou} padá na faktu frameworku, který nástroj měří (konkatenace nad nullable sloupcem, agregát nad celým výsledkem, řetěz množinových operací do EclipseLinku, \texttt{CASE} v~\texttt{ELSE}); a žádná případová studie nedojde ke čtvrtému stupni, protože model nemůže znát schéma projektu, které nástroj bere z~katalogu. Prompt Cofoundry má 105--113~tisíc tokenů a pokus trvá 10--12 minut; cena 258 pokusů podle ceníku je 111,53~USD, zatímco nástroj je nad každou buňkou pod sekundou a bez ceny.
\item \textbf{Co je výsledek protokolu a co modelu.} Protokolu: shoda tvaru souboru se soudcem a oprava obálky; modelu: obsah. Rozptyl tří pokusů je při teplotě 0 nulový (Qwen) a při nenastavitelné teplotě výrobce malý (jedna buňka ze třinácti u~Fable, oběma směry u~Opus); tři pokusy zůstávají.
\end{itemize}

Varianty \texttt{schema} (zero-shot s~DDL katalogu, měřící vliv databázových metadat podle T6), \texttt{rag}, \texttt{repair} a \texttt{agent} jsou v~protokolu a v~manifestu připravené, ale nezaznamenané: patří navazující diplomové práci.

\subsection{Reprodukce pro recenzenta}\label{sec:reprodukce}

Reprodukční návod je soubor \path{docs/reproduction.md} v~repozitáři (česky), psaný pro čtenáře, který má stroj s~Dockerem a chce ověřit, co zpráva tvrdí: v~pořadí obě sady v~kontejnerech, experimentální běh, běh FIBEN, srovnání s~jazykovým modelem včetně nasazení Qwenu, a volitelně soudce LDBC. U~každého kroku říká, co se spouští, kde je vidět postup (log kontejneru, tabulka \texttt{attempts.tsv} psaná po pokusech, soubory běhu psané po dotazech), jak má vypadat výsledek (náhled souhrnu testů, tvar \texttt{summary.md}, řádek záznamu) a jaké chyby mohou nastat a co znamenají. Návod odkazuje na \path{ORMConvertor/README.md} pro úplný výčet proměnných a na \texttt{docs/runs/} pro čísla, s~nimiž se výsledek srovnává; nic v~něm není číslo, které by příští změna posunula.
% ===== 06-zaver.local.tex =====
\section{Závěr}\label{sec:zaver}

Projekt splnil, co specifikace nazvala jádrem rozšíření: společná mezireprezentace a překladový přístup unesly druhý programovací jazyk a druhý technologický ekosystém. Tři javovské frameworky vstoupily jako tenké wrappery v~C\# bez zásahu do orchestrace (S1), překlad jde ve všech 36 směrech šesti frameworků, přes hranici ekosystémů oběma směry, a to, co každý směr dělá s~každou konstrukcí, je vymezené rozhodnutím a měřené na čtyřech stupních v~zafixovaném prostředí. Tabulka~\ref{tab:pozadavky} shrnuje stav požadavků podle §9 \texttt{architecture.md} a sekce \emph{Guarantees}, kde je kanonické znění.

\begin{table}[ht]
\centering\small
\begin{tabularx}{\textwidth}{L{2.4cm} L{3.2cm} Y}
\toprule
\textbf{Požadavky} & \textbf{Stav} & \textbf{Poznámka} \\
\midrule
F1--F6 & nárokované & složené klíče, cizí klíče na ně, katalog, slučování zdrojů, úplné mapování z~neúplného vstupu \\
F7, F8, F9 & nárokované v~užším rozsahu & tři frameworky, ne ekosystém; dynamický příkaz MyBatisu a líné načtení EclipseLinku vyňaté; stupně ověření v~kontejneru \\
F10, F12 & nárokované & překlad přes hranici v~každé ze čtyř uspořádaných dvojic ekosystémů; javovská sada, která si počítá velikost \\
F11, F13, F14 & nárokované v~užším rozsahu & syntaktickou správnost dokazují testy; měřítko porovnání pevné; bez zobrazení mezireprezentace a bez procházení repozitáře \\
F15 & vyňatý & Advisor měří jen Dapper a EF~Core (rozh.~171) \\
S1--S3 & nárokované & modularita, determinismus, výkon překladu měřený (0,2--0,7~s proti mezi 30~s) \\
S4 & 2.~věta nárokovaná & žádné přihlašovací údaje v~artefaktech ani textu o~běhu; izolace spouštění vyňata s~Advisorem \\
S5, S6, S7 & nárokované v~užším rozsahu & kontejnery pro systém, databázi a obě sady; záznam převodu bez výsledků kompilace; validace na klientovi jako pomocník \\
T1, T2, T3 & nárokované v~užším rozsahu & vstup případové studie, matice kategorií, měřidlo experimentálního běhu; čísla jsou záznam verze \\
T4, T5, T6 & vyňaté & protokol a infrastruktura stojí, čtyři záznamy srovnání; plné experimenty v~diplomové práci \\
T7 & vyňatý & s~Advisorem \\
\bottomrule
\end{tabularx}
\caption{Stav požadavků ve vydání 2.3.0.}
\label{tab:pozadavky}
\end{table}

Výzkumný přínos je dvojí a oba díly jsou doložené záznamy, ne tvrzením. Za prvé, \emph{cross-language použitelnost} přístupu: týž model a tytéž překladové kroky dávají artefakty, které \texttt{javac} přeloží, Hibernate, EclipseLink a MyBatis přijmou a které nad SQL Serverem vracejí tytéž řádky jako zdroj, nad maticí 59 kategorií i nad 237 dotazy benchmarku FIBEN, který nástroj nepsal; validační sada LDBC, jediná nezávislá pravda mezi soudci, souhlasí s~každým čtením ve všech šesti frameworcích. Za druhé, \emph{deterministická baseline} pro srovnání s~jazykovými modely a společné měřidlo: tatáž cesta soudí artefakt nástroje i modelu a první čtyři záznamy ukazují, kde je rozdíl -- ne v~jazyce dotazu, který silný model píše správně, ale ve faktech, která model nezná (schéma projektu, jméno sestavení, chování připnuté verze frameworku) a nástroj bere z~katalogu a z~rozhodnutí podložených měřením.

Co zbývá, nese \texttt{docs/open-items.md} a je trojího druhu. Pro odevzdání: scénář živé ukázky u~obhajoby nad čerstvě sestavenými obrazy vydání a kontrola, že veřejný repozitář vede čtenáře od spuštění přes nárok k~důkazům. Pro navazující diplomovou práci: plné experimenty T4--T6 -- varianta \texttt{schema} jako měření vlivu databázových metadat, RAG s~dokumentací cílového frameworku, opravný cyklus nad verdiktem soudce a agentní varianta, ablace --, první dokončený běh SEDE a jeho soudce, výňatek Cofoundry dovedený do čtvrtého stupně. Pro příští MAJOR vydání \texttt{3.0.0}: Advisor nad všemi šesti frameworky s~javovskou větví měření (rozh.~146, 147), srovnání s~jednotlivými frameworky a jednoduchou heuristikou (T7) a izolace spouštění cizího kódu (S4). Výsledky projektu -- přístup, měřicí infrastruktura a čísla záznamů -- jsou podkladem pro publikaci, kterou specifikace jmenuje jako využití výsledků.

\begin{thebibliography}{99}

\bibitem{spec}
ŠTOCHEL, Miroslav. \emph{Rozšíření nástroje ORMorpher: Specification of team software project}. Schválená specifikace výzkumného projektu, MFF UK, 2026. V~repozitáři projektu jako \path{docs/specifikace.tex}.

\bibitem{zamer}
ŠTOCHEL, Miroslav. \emph{Rozšíření nástroje ORMorpher: záměr projektu}. Schválený záměr, MFF UK, 2026. V~repozitáři projektu jako \texttt{docs/zamer.tex}.

\bibitem{fork-repo}
ŠTOCHEL, Miroslav. \emph{orm-convertor} (ORMConvertor). Repozitář projektu, vydání 2.3.0. Dostupné z: \url{https://github.com/MiraStochel/orm-convertor}

\bibitem{abraham-repo}
ABRAHÁM, Milan. \emph{Framework-Agnostic Query Adaptation: Ensuring SQL Compatibility Across .NET Database Frameworks}. Diplomová práce, Univerzita Karlova, Praha, 2025. Repozitář prototypu: \url{https://github.com/milan252525/orm-convertor}

\bibitem{ormorpher-ase}
ABRAHÁM, Milan; KOUPIL, Pavel. ORMorpher: An Interactive Framework for ORM Translation and Optimization. In: \emph{40th IEEE/ACM International Conference on Automated Software Engineering (ASE 2025)}. Seoul, 2025.

\bibitem{ormorpher-jss}
ABRAHÁM, Milan; KOUPIL, Pavel. \emph{A Unified Framework for Object-Relational Mapping Translation and Performance-Aware Selection}. Rozšířená časopisecká verze, preprint, 2026.

\bibitem{adam-web}
ADAPTIVE DATA MANAGEMENT RESEARCH GROUP. Dostupné z: \url{https://www.cnaip.cz/en/entity/adaptive-data-management-research-group}

\bibitem{mmcat-intro}
MM-cat Docs. \emph{Motivation}. Dostupné z: \url{https://mmcatdb.com/getting-started/introduction/} [čteno 9.~11.~2026]

\bibitem{mmcat-aimm}
MM-cat Docs. \emph{Database Mapping Optimization}. Dostupné z: \url{https://mmcatdb.com/aimm/} [čteno 9.~11.~2026]

\bibitem{koupil-holubova-2022}
KOUPIL, Pavel; HOLUBOVÁ, Irena. A unified representation and transformation of multi-model data using category theory. \emph{Journal of Big Data}, 2022, 9, 61. \url{https://doi.org/10.1186/s40537-022-00613-3}

\bibitem{mm-quecat}
KOUPIL, Pavel; CRHA, Daniel; HOLUBOVÁ, Irena. MM-quecat: A Tool for Unified Querying of Multi-Model Data. In: \emph{EDBT 2023}. \url{https://doi.org/10.48786/edbt.2023.76}

\bibitem{mm-evoque}
KOUPIL, Pavel; BÁRTÍK, Jáchym; HOLUBOVÁ, Irena. MM-evoque: Query Synchronisation in Multi-Model Databases. In: \emph{EDBT 2024}. \url{https://doi.org/10.48786/edbt.2024.78}

\bibitem{mm-evocat}
KOUPIL, Pavel; BÁRTÍK, Jáchym; HOLUBOVÁ, Irena. MM-evocat: A Tool for Modelling and Evolution Management of Multi-Model Data. In: \emph{CIKM 2022}. \url{https://doi.org/10.1145/3511808.3557180}

\bibitem{mm-mapsearch}
KOUPIL, Pavel et al. \emph{mm-mapsearch}. Stránka nástroje. Dostupné z: \url{https://www.ksi.mff.cuni.cz/~koupil/software/mm-mapsearch/}

\bibitem{ldbc}
LDBC. \emph{Social Network Benchmark, Interactive workload v1}. Datové sady a validační sady: \url{https://datasets.ldbcouncil.org/interactive-v1/}

\bibitem{fiben}
IBM Research. \emph{FIBEN: Financial Benchmark}. Repozitář: \url{https://github.com/IBM/fiben-benchmark}

\bibitem{sede}
Stack Exchange Data Explorer query set (SEDE). Repozitář: \url{https://github.com/hirupert/sede}; data: výpis dba.stackexchange.com, Internet Archive, položka \texttt{stackexchange}.

\bibitem{qwen}
Qwen Team. \emph{Qwen3.5-9B}. Hugging Face; kvantizace Q4\_K\_M: \path{unsloth/Qwen3.5-9B-GGUF}.

\bibitem{llamacpp}
GERGANOV, Georgi et al. \emph{llama.cpp}. Serverový obraz \path{ghcr.io/ggml-org/llama.cpp:server}. \url{https://github.com/ggml-org/llama.cpp}

\end{thebibliography}

\end{document}
